\chapter{Asymptotic Normality}\label{ch:asymptotic-normality}

In the classical linear model with fixed regressors and normally distributed
i.i.d.\ errors, the least squares estimator $\hat{\bfbeta}_n$ is distributed as a multivariate
normal with $E(\hat{\bfbeta}_n) = \bfbeta_o$ and $\mathrm{var}(\hat{\bfbeta}_n) = \sigma_o^2(\mathbf{X}'\mathbf{X})^{-1}$ for any sample size
$n$. This fact forms the basis for statistical tests of hypotheses, based typically on $t$- and $F$-statistics. When the sample size is large, econometric
estimators such as $\hat{\bfbeta}_n$ have a distribution that is approximately normal
under much more general conditions, and this fact forms the basis for large
sample statistical tests of hypotheses. In this chapter we study the tools
used in determining the asymptotic distribution of $\hat{\bfbeta}_n$, how this asymptotic distribution can be used to test hypotheses in large samples, and how
asymptotic efficiency can be obtained.

\section{Convergence in Distribution}\label{sec:convergence-distribution}

The most fundamental concept is that of convergence in distribution.

\begin{definition}\label{def:4.1}
Let $\{\mathbf{b}_n\}$ be a sequence of random finite-dimensional vectors with joint distribution functions $\{F_n\}$. If $F_n(\mathbf{z}) \to F(\mathbf{z})$ as $n \to \infty$ for every continuity point $\mathbf{z}$, where $F$ is the distribution function of a random variable $\mathcal{Z}$, then $\mathbf{b}_n$ converges in distribution to the random variable $\mathcal{Z}$, denoted $\mathbf{b}_n \convd \mathcal{Z}$.
\end{definition}

Heuristically, the distribution of $\mathbf{b}_n$ gets closer and closer to that of the
random variable $\mathcal{Z}$, so the distribution $F$ can be used as an approximation
to the distribution of $\mathbf{b}_n$. When $\mathbf{b}_n \convd \mathcal{Z}$, we also say that $\mathbf{b}_n$ \emph{converges in
law} to $\mathcal{Z}$ (written $\mathbf{b}_n \xrightarrow{L} \mathcal{Z}$), or that $\mathbf{b}_n$ is \emph{asymptotically distributed} as
$F$, denoted $\mathbf{b}_n \overset{A}{\sim} F$. Then $F$ is called the \emph{limiting distribution} of $\mathbf{b}_n$. Note
that the convergence specified by this definition is pointwise and only has
to occur at points $\mathbf{z}$ where $F$ is continuous (the ``continuity points'').

\begin{example}\label{ex:4.2}
Let $\{\mathbf{b}_n\}$ be a sequence of i.i.d.\ random variables with distribution function $F$. Then (trivially) $F$ is the limiting distribution of $\mathbf{b}_n$.
\end{example}

This illustrates the fact that convergence in distribution is a very weak
convergence concept and by itself implies nothing about the convergence
of the sequence of random variables.

\begin{example}\label{ex:4.3}
Let $\{\mathcal{Z}_t\}$ be i.i.d.\ random variables with mean $\mu$ and finite variance $\sigma^2 > 0$. Define
\[
b_n \equiv \frac{\bar{\mathcal{Z}}_n - E(\bar{\mathcal{Z}}_n)}{(\mathrm{var}(\bar{\mathcal{Z}}_n))^{1/2}} = n^{-1/2} \sum_{t=1}^{n} (\mathcal{Z}_t - \mu)/\sigma.
\]
Then by the Lindeberg-L\'evy central limit theorem (Theorem~5.2), $b_n \overset{A}{\sim} N(0,1)$.
\end{example}

In other words, the sample mean $\bar{\mathcal{Z}}_n$ of i.i.d.\ observations, when standardized, has a distribution that approaches the standard normal distribution.
This result actually holds under rather general conditions on the sequence
$\{\mathcal{Z}_t\}$. The conditions under which this convergence occurs are studied at
length in the next chapter. In this chapter, we simply assume that such
conditions are satisfied, so convergence in distribution is guaranteed.

Convergence in distribution is meaningful even when the limiting distribution is that of a degenerate random variable.

\begin{lemma}\label{lem:4.4}
Suppose $b_n \convp b$ (a constant). Then $b_n \overset{A}{\sim} F_b$, where $F_b$ is the distribution function of a random variable $\mathcal{Z}$ that takes the value $b$ with probability one (i.e., $b_n \convd b$). Also, if $b_n \overset{A}{\sim} F_b$, then $b_n \convp b$.
\end{lemma}

\begin{proof}
See Rao (1973, p.~120).
\end{proof}

In other words, convergence in probability to a constant implies convergence
in distribution to that constant. The converse is also true.

A useful implication of convergence in distribution is the following lemma.

\begin{lemma}\label{lem:4.5}
If $b_n \convd \mathcal{Z}$, then $b_n$ is $O_p(1)$.
\end{lemma}

\begin{proof}
Recall that $b_n$ is $O_p(1)$ if, given any $\delta > 0$, $P[|b_n| > \Delta_\delta] < \delta$
for some $\Delta_\delta < \infty$ and all $n > N_\delta$. Because $b_n \to \mathcal{Z}$, $P[|b_n| > \Delta_\delta] \to
P[|\mathcal{Z}| > \Delta_\delta]$, provided (without loss of generality) that $\Delta_\delta$ and $-\Delta_\delta$ are
continuity points of the distribution of $\mathcal{Z}$. Thus $|P[|b_n| > \Delta_\delta] - P[|\mathcal{Z}| >
\Delta_\delta]| < \delta$ for all $n > N_\delta$, so that $P[|b_n| > \Delta_\delta] < \delta + P[|\mathcal{Z}| > \Delta_\delta]$ for
all $n > N_\delta$. Because $P[|\mathcal{Z}| > \Delta_\delta] < \delta$ for $\Delta_\delta$ sufficiently large, we have
$P[|b_n| > \Delta_\delta] < 2\delta$ for $\Delta_\delta$ sufficiently large and all $n > N_\delta$.
\end{proof}

This allows us to establish the next useful lemma.

\begin{lemma}[Product rule]\label{lem:4.6}
Recall from Corollary~2.36 that if $\mathbf{A}_n = o_p(1)$ and $\mathbf{b}_n = O_p(1)$, then $\mathbf{A}_n \mathbf{b}_n = o_p(1)$. Hence, if $\mathbf{A}_n \convp \mathbf{0}$ and $\mathbf{b}_n \convd \mathcal{Z}$, then $\mathbf{A}_n \mathbf{b}_n \convp \mathbf{0}$.
\end{lemma}

In turn, this result is often used in conjunction with the following result,
which is one of the most useful of those relating convergence in probability
and convergence in distribution.

\begin{lemma}[Asymptotic equivalence]\label{lem:4.7}
Let $\{\mathbf{a}_n\}$ and $\{\mathbf{b}_n\}$ be two sequences of random vectors. If $\mathbf{a}_n - \mathbf{b}_n \convp \mathbf{0}$ and $\mathbf{b}_n \convd \mathcal{Z}$, then $\mathbf{a}_n \convd \mathcal{Z}$.
\end{lemma}

\begin{proof}
See Rao (1973, p.~123).
\end{proof}

This result is helpful in situations in which we wish to find the asymptotic
distribution of $\mathbf{a}_n$ but cannot easily do so directly. Often, however, it is easy
to find a $\mathbf{b}_n$ that has a known asymptotic distribution and that satisfies
$\mathbf{a}_n - \mathbf{b}_n \convp \mathbf{0}$. Lemma~4.7 then ensures that $\mathbf{a}_n$ has the same limiting
distribution as $\mathbf{b}_n$ and we say that $\mathbf{a}_n$ is ``asymptotically equivalent'' to
$\mathbf{b}_n$. The joint use of Lemmas~4.6 and 4.7 is the key to the proof of the
asymptotic normality results for the OLS and IV estimators.

Another useful tool in the study of convergence in distribution is the
characteristic function.

\begin{definition}\label{def:4.8}
Let $\mathcal{Z}$ be a $k \times 1$ random vector with distribution function $F$. The characteristic function of $\mathcal{Z}$ is defined as $f(\boldsymbol{\lambda}) \equiv E(\exp(i\boldsymbol{\lambda}'\mathcal{Z}))$, where $i^2 = -1$ and $\boldsymbol{\lambda}$ is a $k \times 1$ real vector.
\end{definition}

\begin{example}\label{ex:4.9}
Let $\mathcal{Z}$ be a nonstochastic real number, $\mathcal{Z} = c$. Then $f(\lambda) = E(\exp(i\lambda \mathcal{Z})) = E(\exp(i\lambda c)) = \exp(i\lambda c)$.
\end{example}

\begin{example}\label{ex:4.10}
$(i)$ Let $\mathcal{Z} \sim N(\mu, \sigma^2)$. Then $f(\lambda) = \exp(i\lambda\mu - \lambda^2\sigma^2/2)$.
$(ii)$ Let $\mathcal{Z} \sim N(\boldsymbol{\mu}, \boldsymbol{\Sigma})$, where $\boldsymbol{\mu}$ is $k \times 1$ and $\boldsymbol{\Sigma}$ is $k \times k$. Then $f(\boldsymbol{\lambda}) = \exp(i\boldsymbol{\lambda}'\boldsymbol{\mu} - \boldsymbol{\lambda}'\boldsymbol{\Sigma}\boldsymbol{\lambda}/2)$.
\end{example}

A useful table of characteristic functions is given by Lukacs (1970, p.~18).

Because the characteristic function is the Fourier transformation of the
probability density function, it has the property that any characteristic
function uniquely determines a distribution function, as formally expressed
by the next result.

\begin{theorem}[Uniqueness theorem]\label{thm:4.11}
Two distribution functions are identical if and only if their characteristic functions are identical.
\end{theorem}

\begin{proof}
See Lukacs (1974, p.~14).
\end{proof}

Thus the behavior of a random variable can be studied either through
its distribution function or its characteristic function, whichever is more
convenient.

\begin{example}\label{ex:4.12}
The distribution of a linear transformation of a random
variable is easily found using the characteristic function. Consider $\mathcal{Y} = \mathbf{A}\mathcal{Z}$, where $\mathbf{A}$ is a $q \times k$ matrix and $\mathcal{Z}$ is a random $k$-vector. Let $\boldsymbol{\theta}$ be $q \times 1$.
Then
\begin{eqnarray*}
f_{\mathcal{Y}}(\boldsymbol{\theta}) &=& E(\exp(i\boldsymbol{\theta}'\mathcal{Y})) = E(\exp(i\boldsymbol{\theta}'\mathbf{A}\mathcal{Z}))\\
&=& E(\exp(i\boldsymbol{\lambda}'\mathcal{Z})) = f_{\mathcal{Z}}(\boldsymbol{\lambda}),
\end{eqnarray*}
defining $\boldsymbol{\lambda} = \mathbf{A}'\boldsymbol{\theta}$. Hence if $\mathcal{Z} \sim N(\boldsymbol{\mu}, \boldsymbol{\Sigma})$,
\begin{eqnarray*}
f_{\mathcal{Y}}(\boldsymbol{\theta}) &=& f_{\mathcal{Z}}(\boldsymbol{\lambda}) = \exp(i\boldsymbol{\lambda}'\boldsymbol{\mu} - \boldsymbol{\lambda}'\boldsymbol{\Sigma}\boldsymbol{\lambda}/2)\\
&=& \exp(i\boldsymbol{\theta}'\mathbf{A}\boldsymbol{\mu} - \boldsymbol{\theta}'\mathbf{A}\boldsymbol{\Sigma}\mathbf{A}'\boldsymbol{\theta}/2),
\end{eqnarray*}
so that $\mathcal{Y} \sim N(\mathbf{A}\boldsymbol{\mu}, \mathbf{A}\boldsymbol{\Sigma}\mathbf{A}')$ by the uniqueness theorem.
\end{example}

Other useful facts regarding characteristic functions are the following.

\begin{proposition}\label{prop:4.13}
Let $\mathcal{Y} = a\mathcal{Z} + b$, $a, b \in \mathbb{R}$. Then
\[
f_{\mathcal{Y}}(\lambda) = f_{\mathcal{Z}}(a\lambda)\exp(i\lambda b).
\]
\end{proposition}

\begin{proof}
\begin{eqnarray*}
f_{\mathcal{Y}}(\lambda) &=& E(\exp(i\lambda\mathcal{Y})) = E(\exp(i\lambda(a\mathcal{Z} + b)))\\
&=& E(\exp(i\lambda a\mathcal{Z})\exp(i\lambda b))\\
&=& E(\exp(i\lambda a\mathcal{Z}))\exp(i\lambda b) = f_{\mathcal{Z}}(\lambda a)\exp(i\lambda b).
\end{eqnarray*}
\end{proof}

\begin{proposition}\label{prop:4.14}
Let $\mathcal{Y}$ and $\mathcal{Z}$ be independent. Then if $\mathcal{X} = \mathcal{Y} + \mathcal{Z}$, $f_{\mathcal{X}}(\lambda) = f_{\mathcal{Y}}(\lambda)f_{\mathcal{Z}}(\lambda)$.
\end{proposition}

\begin{proof}
\begin{eqnarray*}
f_{\mathcal{X}}(\lambda) &=& E(\exp(i\lambda\mathcal{X})) = E(\exp(i\lambda(\mathcal{Y} + \mathcal{Z})))\\
&=& E(\exp(i\lambda\mathcal{Y})\exp(i\lambda\mathcal{Z}))\\
&=& E(\exp(i\lambda\mathcal{Y}))E(\exp(i\lambda\mathcal{Z}))
\end{eqnarray*}
by independence. Hence $f_{\mathcal{X}}(\lambda) = f_{\mathcal{Y}}(\lambda)f_{\mathcal{Z}}(\lambda)$.
\end{proof}

\begin{proposition}\label{prop:4.15}
If the $k$th moment $\mu_k$ of a distribution function $F$ exists, then the characteristic function $f$ of $F$ can be differentiated $k$ times and $f^{(k)}(0) = i^k\mu_k$, where $f^{(k)}$ is the $k$th derivative of $f$.
\end{proposition}

\begin{proof}
This is an immediate corollary of Lukacs (1970, Corollary~3 to
Theorem~2.3.1, p.~22).
\end{proof}

\begin{example}\label{ex:4.16}
Suppose that $\mathcal{Z} \sim N(0, \sigma^2)$. Then $f'(0) = 0$, $f''(0) = -\sigma^2$, $f'''(0) = 0$, etc.
\end{example}

The main result of use in studying convergence in distribution is the
following.

\begin{theorem}[Continuity theorem]\label{thm:4.17}
Let $\{\mathbf{b}_n\}$ be a sequence of random $k \times 1$ vectors with characteristic functions $\{f_n(\boldsymbol{\lambda})\}$. If $\mathbf{b}_n \convd \mathcal{Z}$, then for every $\boldsymbol{\lambda}$, $f_n(\boldsymbol{\lambda}) \to f(\boldsymbol{\lambda})$, where $f(\boldsymbol{\lambda}) = E(\exp(i\boldsymbol{\lambda}'\mathcal{Z}))$. Further, if for every $\boldsymbol{\lambda}$, $f_n(\boldsymbol{\lambda}) \to f(\boldsymbol{\lambda})$ and $f$ is continuous at $\boldsymbol{\lambda} = \mathbf{0}$, then $\mathbf{b}_n \convd \mathcal{Z}$, where $f(\boldsymbol{\lambda}) = E(\exp(i\boldsymbol{\lambda}'\mathcal{Z}))$.
\end{theorem}

\begin{proof}
See Lukacs (1970, pp.~49--50).
\end{proof}

This result essentially says that convergence in distribution is equivalent
to convergence of characteristic functions. The usefulness of the result is
that often it is much easier to study the limiting behavior of characteristic functions than distribution functions. If the sequence of characteristic
functions $f_n$ converges to a function $f$ that is continuous at $\boldsymbol{\lambda} = \mathbf{0}$, this
theorem guarantees that the function $f$ is a characteristic function and that
the limiting distribution $F$ of $\mathbf{b}_n$ is that corresponding to the characteristic
function $f(\boldsymbol{\lambda})$.

In all the cases that follow, the limiting distribution $F$ will be either that
of a degenerate random variable (following from convergence in probability
to a constant) or a multivariate normal distribution (following from an
appropriate central limit theorem). In the latter case, it is often convenient
to standardize the random variables so that the asymptotic distribution is
unit multivariate normal. To do this we can use the matrix square root.

\begin{exercises}
\item\label{ex:4.18} \textit{Prove the following. Let\/ $\mathbf{V}$ be a positive (semi) definite
symmetric matrix. Then there exists a positive (semi) definite symmetric
matrix square root $\mathbf{V}^{1/2}$ such that the elements of\/ $\mathbf{V}^{1/2}$ are continuous
functions of\/ $\mathbf{V}$ and $\mathbf{V}^{1/2}\mathbf{V}^{1/2} = \mathbf{V}$. (Hint: Express $\mathbf{V}$ as $\mathbf{V} = \mathbf{Q}'\mathbf{D}\mathbf{Q}$, where
$\mathbf{Q}$ is an orthogonal matrix and $\mathbf{D}$ is diagonal with the eigenvalues of\/ $\mathbf{V}$
along the diagonal.)}

\item\label{ex:4.19} \textit{Show that if $\mathcal{Z} \sim N(\mathbf{0}, \mathbf{V})$, then $\mathbf{V}^{-1/2}\mathcal{Z} \sim N(\mathbf{0}, \mathbf{I})$, provided $\mathbf{V}$ is positive definite, where $\mathbf{V}^{-1/2} = (\mathbf{V}^{1/2})^{-1}$.}
\end{exercises}

\begin{definition}\label{def:4.20}
Let $\{\mathbf{b}_n\}$ be a sequence of random vectors. If there exists a sequence of matrices $\{\mathbf{V}_n\}$ such that $\mathbf{V}_n$ is nonsingular for all $n$ sufficiently large and $\mathbf{V}_n^{-1/2}\mathbf{b}_n \overset{A}{\sim} N(\mathbf{0}, \mathbf{I})$, then $\mathbf{V}_n$ is called the asymptotic covariance matrix of $\mathbf{b}_n$, denoted $\mathrm{avar}(\mathbf{b}_n)$.
\end{definition}

When $\mathrm{var}(\mathbf{b}_n)$ is finite, we can usually define $\mathbf{V}_n = \mathrm{var}(\mathbf{b}_n)$. Note that
the behavior of $\mathbf{b}_n$ is not restricted to require that $\mathbf{V}_n$ converge to any
limit, although it may. Generally, however, we will at least require that the
smallest eigenvalues of $\mathbf{V}_n$ and $\mathbf{V}_n^{-1}$ are uniformly bounded away from zero
for all $n$ sufficiently large. Even when $\mathrm{var}(\mathbf{b}_n)$ is not finite, the asymptotic
covariance matrix can exist, although in such cases we cannot set $\mathbf{V}_n =
\mathrm{var}(\mathbf{b}_n)$.

\begin{example}\label{ex:4.21}
Define $b_n = \mathcal{Z} + \mathcal{Y}/n$, where $\mathcal{Z} \sim N(0,1)$ and $\mathcal{Y}$ is Cauchy, independent of $\mathcal{Z}$. Then $\mathrm{var}(b_n)$ is infinite for every $n$, but $b_n \overset{A}{\sim} N(0,1)$ as a consequence of Lemma~4.7. Hence $\mathrm{avar}(b_n) = 1$.
\end{example}

Given a sequence $\{\mathbf{V}_n^{-1/2}\mathbf{b}_n\}$ that converges in distribution, we shall often be interested in the behavior of linear combinations of $\mathbf{b}_n$, say, $\{\mathbf{A}_n\mathbf{b}_n\}$,
where $\mathbf{A}_n$, like $\mathbf{V}_n^{-1/2}$, is not required to converge to a particular limit. We
can use characteristic functions to study the behavior of these sequences
by making use of the following corollary to the continuity theorem.

\begin{corollary}\label{cor:4.22}
If $\boldsymbol{\lambda} \in \mathbb{R}^k$ and a sequence $\{f_n(\boldsymbol{\lambda})\}$ of characteristic functions converges to a characteristic function $f(\boldsymbol{\lambda})$, then the convergence is uniform in every compact subset of $\mathbb{R}^k$.
\end{corollary}

\begin{proof}
This is a straightforward extension of Lukacs (1970, p.~50).
\end{proof}

This result says that in any compact subset of $\mathbb{R}^k$ the distance between
$f_n(\boldsymbol{\lambda})$ and $f(\boldsymbol{\lambda})$ does not depend on $\boldsymbol{\lambda}$, but only on $n$. This fact is crucial
to establishing the next result.

\begin{lemma}\label{lem:4.23}
Let $\{\mathbf{b}_n\}$ be a sequence of random $k \times 1$ vectors with characteristic functions $\{f_n(\boldsymbol{\lambda})\}$, and suppose $f_n(\boldsymbol{\lambda}) \to f(\boldsymbol{\lambda})$. If $\{\mathbf{A}_n\}$ is any sequence of $q \times k$ nonstochastic matrices such that $\mathbf{A}_n = O(1)$, then the sequence $\{\mathbf{A}_n\mathbf{b}_n\}$ has characteristic functions $\{f_n^*(\boldsymbol{\theta})\}$, where $\boldsymbol{\theta}$ is $q \times 1$, such that for every $\boldsymbol{\theta}$, $f_n^*(\boldsymbol{\theta}) - f(\mathbf{A}_n'\boldsymbol{\theta}) \to 0$.
\end{lemma}

\begin{proof}
From Example~4.12, $f_n^*(\boldsymbol{\theta}) = f_n(\mathbf{A}_n'\boldsymbol{\theta})$. For fixed $\boldsymbol{\theta}$, $\boldsymbol{\lambda}_n \equiv \mathbf{A}_n'\boldsymbol{\theta}$ takes
values in a compact region of $\mathbb{R}^k$, say, $\mathcal{N}_\theta$, for all $n$ sufficiently large because
$\mathbf{A}_n = O(1)$. Because $f_n(\boldsymbol{\lambda}) \to f(\boldsymbol{\lambda})$, we have $f_n(\boldsymbol{\lambda}_n) - f(\boldsymbol{\lambda}_n) \to 0$ uniformly
for all $\boldsymbol{\lambda}_n$ in $\mathcal{N}_\theta$, by Corollary~4.22. Hence for fixed $\boldsymbol{\theta}$, $f_n(\mathbf{A}_n'\boldsymbol{\theta}) - f(\mathbf{A}_n'\boldsymbol{\theta}) =
f_n^*(\boldsymbol{\theta}) - f(\mathbf{A}_n'\boldsymbol{\theta}) \to 0$ for any $O(1)$ sequence $\{\mathbf{A}_n\}$. Because $\boldsymbol{\theta}$ is arbitrary,
the result follows.
\end{proof}

The following consequence of this result is used many times below.

\begin{corollary}\label{cor:4.24}
Let $\{\mathbf{b}_n\}$ be a sequence of random $k \times 1$ vectors such that $\mathbf{V}_n^{-1/2}\mathbf{b}_n \overset{A}{\sim} N(\mathbf{0}, \mathbf{I})$, where $\{\mathbf{V}_n\}$ and $\{\mathbf{V}_n^{-1}\}$ are $O(1)$. Let $\{\mathbf{A}_n\}$ be a $O(1)$ sequence of (nonstochastic) $q \times k$ matrices with full row rank $q$ for all $n$ sufficiently large, uniformly in $n$. Then the sequence $\{\mathbf{A}_n\mathbf{b}_n\}$ is such that $\boldsymbol{\Gamma}_n^{-1/2}\mathbf{A}_n\mathbf{b}_n \overset{A}{\sim} N(\mathbf{0}, \mathbf{I})$, where $\boldsymbol{\Gamma}_n \equiv \mathbf{A}_n\mathbf{V}_n\mathbf{A}_n'$ and $\boldsymbol{\Gamma}_n$ and $\boldsymbol{\Gamma}_n^{-1}$ are $O(1)$.
\end{corollary}

\begin{proof}
$\boldsymbol{\Gamma}_n = O(1)$ by Lemma~2.19. $\boldsymbol{\Gamma}_n^{-1} = O(1)$ because $\boldsymbol{\Gamma}_n = O(1)$
and $\det(\boldsymbol{\Gamma}_n) > \delta > 0$ for all $n$ sufficiently large, given the conditions on
$\{\mathbf{A}_n\}$ and $\{\mathbf{V}_n\}$. Let $f_n^*(\boldsymbol{\theta})$ be the characteristic function of $\boldsymbol{\Gamma}_n^{-1/2}\mathbf{A}_n\mathbf{b}_n =$
$\boldsymbol{\Gamma}_n^{-1/2}\mathbf{A}_n\mathbf{V}_n^{1/2}\mathbf{V}_n^{-1/2}\mathbf{b}_n$. Because $\boldsymbol{\Gamma}_n^{-1/2}\mathbf{A}_n\mathbf{V}_n^{1/2} = O(1)$, Lemma~4.23 applies, implying $f_n^*(\boldsymbol{\theta}) - f(\mathbf{V}_n^{1/2}\mathbf{A}_n'\boldsymbol{\Gamma}_n^{-1/2}\boldsymbol{\theta}) \to 0$, where $f(\boldsymbol{\lambda}) = \exp(-\boldsymbol{\lambda}'\boldsymbol{\lambda}/2)$,
the limiting characteristic function of $\mathbf{V}_n^{-1/2}\mathbf{b}_n$. Now $f(\mathbf{V}_n^{1/2}\mathbf{A}_n'\boldsymbol{\Gamma}_n^{-1/2}\boldsymbol{\theta}) =
\exp(-\boldsymbol{\theta}'\boldsymbol{\Gamma}_n^{-1/2}\mathbf{A}_n\mathbf{V}_n\mathbf{A}_n'\boldsymbol{\Gamma}_n^{-1/2}\boldsymbol{\theta}/2) = \exp(-\boldsymbol{\theta}'\boldsymbol{\theta}/2)$ by definition of $\boldsymbol{\Gamma}_n^{-1/2}$.
Hence $f_n^*(\boldsymbol{\lambda}) - \exp(-\boldsymbol{\theta}'\boldsymbol{\theta}/2) \to 0$, so $\boldsymbol{\Gamma}_n^{-1/2}\mathbf{A}_n\mathbf{b}_n \overset{A}{\sim} N(\mathbf{0}, \mathbf{I})$ by the continuity theorem (Theorem~4.17).
\end{proof}

This result allows us to complete the proof of the following general
asymptotic normality result for the least squares estimator.

\begin{theorem}\label{thm:4.25}
\emph{Given}
\begin{enumerate}[label=(\roman*)]
\item $\mathbf{Y}_t = \mathbf{X}_t'\bfbeta_o + \eps_t$, \quad $t = 1, 2, \ldots$\,, $\bfbeta_o \in \mathbb{R}^k$;
\item $\mathbf{V}_n^{-1/2}n^{-1/2}\mathbf{X}'\bfeps \overset{A}{\sim} N(\mathbf{0}, \mathbf{I})$, where $\mathbf{V}_n \equiv \mathrm{var}(n^{-1/2}\mathbf{X}'\bfeps)$ is $O(1)$ and uniformly positive definite;
\item $\mathbf{X}'\mathbf{X}/n - \mathbf{M}_n \convp \mathbf{0}$, where $\mathbf{M}_n = E(\mathbf{X}'\mathbf{X}/n)$ is $O(1)$ and uniformly positive definite.
\end{enumerate}
\emph{Then}
\[
\mathbf{D}_n^{-1/2}\sqrt{n}(\hat{\bfbeta}_n - \bfbeta_o) \overset{A}{\sim} N(\mathbf{0}, \mathbf{I}),
\]
where $\mathbf{D}_n \equiv \mathbf{M}_n^{-1}\mathbf{V}_n\mathbf{M}_n^{-1}$ and $\mathbf{D}_n^{-1}$ are $O(1)$. Suppose in addition that
\begin{enumerate}[label=(\roman*)]
\setcounter{enumi}{3}
\item there exists a matrix $\hat{\mathbf{V}}_n$ positive semidefinite and symmetric such that $\hat{\mathbf{V}}_n - \mathbf{V}_n \convp \mathbf{0}$. Then $\hat{\mathbf{D}}_n - \mathbf{D}_n \convp \mathbf{0}$, where
\[
\hat{\mathbf{D}}_n \equiv (\mathbf{X}'\mathbf{X}/n)^{-1}\hat{\mathbf{V}}_n(\mathbf{X}'\mathbf{X}/n)^{-1}.
\]
\end{enumerate}
\end{theorem}

\begin{proof}
Because $\mathbf{X}'\mathbf{X}/n - \mathbf{M}_n \convp \mathbf{0}$ and $\mathbf{M}_n$ is finite and nonsingular by
$(iii)$, $(\mathbf{X}'\mathbf{X}/n)^{-1}$ and $\hat{\bfbeta}_n$ exist in probability. Given $(i)$ and the existence
of $(\mathbf{X}'\mathbf{X}/n)^{-1}$,
\[
\sqrt{n}(\hat{\bfbeta}_n - \bfbeta_o) = (\mathbf{X}'\mathbf{X}/n)^{-1}n^{-1/2}\mathbf{X}'\bfeps.
\]
Hence, given $(ii)$,
\begin{align*}
\sqrt{n}(\hat{\bfbeta}_n - \bfbeta_o) - \mathbf{M}_n^{-1}n^{-1/2}\mathbf{X}'\bfeps
&= [(\mathbf{X}'\mathbf{X}/n)^{-1} - \mathbf{M}_n^{-1}]\mathbf{V}_n^{1/2}\mathbf{V}_n^{-1/2}n^{-1/2}\mathbf{X}'\bfeps,
\end{align*}
or, premultiplying by $\mathbf{D}_n^{-1/2}$,
\begin{align*}
\mathbf{D}_n^{-1/2}\sqrt{n}(\hat{\bfbeta}_n - \bfbeta_o) &- \mathbf{D}_n^{-1/2}\mathbf{M}_n^{-1}n^{-1/2}\mathbf{X}'\bfeps \\
&= \mathbf{D}_n^{-1/2}[(\mathbf{X}'\mathbf{X}/n)^{-1} - \mathbf{M}_n^{-1}]\mathbf{V}_n^{1/2}\mathbf{V}_n^{-1/2}n^{-1/2}\mathbf{X}'\bfeps.
\end{align*}
The desired result will follow by applying the product rule Lemma~4.6
to the line immediately above, and the asymptotic equivalence Lemma
4.7 to the preceding line. Now $\mathbf{V}_n^{-1/2}n^{-1/2}\mathbf{X}'\bfeps \overset{A}{\sim} N(\mathbf{0}, \mathbf{I})$ by $(ii)$; further,
$\mathbf{D}_n^{-1/2}[(\mathbf{X}'\mathbf{X}/n)^{-1} - \mathbf{M}_n^{-1}]\mathbf{V}_n^{1/2}$ is $o_p(1)$ because $\mathbf{D}_n^{-1/2}$ and $\mathbf{V}_n^{1/2}$ are $O(1)$
given $(ii)$ and $(iii)$, and $[(\mathbf{X}'\mathbf{X}/n)^{-1} - \mathbf{M}_n^{-1}]$ is $o_p(1)$ by Proposition~2.30.
Hence by the product rule Lemma~4.6,
\[
\mathbf{D}_n^{-1/2}\sqrt{n}(\hat{\bfbeta}_n - \bfbeta_o) - \mathbf{D}_n^{-1/2}\mathbf{M}_n^{-1}n^{-1/2}\mathbf{X}'\bfeps \convp \mathbf{0}.
\]
By Lemma~4.7, the asymptotic distribution of $\mathbf{D}_n^{-1/2}\sqrt{n}(\hat{\bfbeta}_n - \bfbeta_o)$ is the
same as that of $\mathbf{D}_n^{-1/2}\mathbf{M}_n^{-1}n^{-1/2}\mathbf{X}'\bfeps$. We find the asymptotic distribution
of this random variable by applying Corollary~4.24, which immediately
yields $\mathbf{D}_n^{-1/2}\mathbf{M}_n^{-1}n^{-1/2}\mathbf{X}'\bfeps \overset{A}{\sim} N(\mathbf{0}, \mathbf{I})$.

Because $(ii)$, $(iii)$, and $(iv)$ hold, $\hat{\mathbf{D}}_n - \mathbf{D}_n \convp \mathbf{0}$ is an immediate consequence of Proposition~2.30.
\end{proof}

The structure of this result is very straightforward. Given the linear
data generating process, we require only that $(\mathbf{X}'\mathbf{X}/n)$ and $(\mathbf{X}'\mathbf{X}/n)^{-1}$ are
$O_p(1)$ and that $n^{-1/2}\mathbf{X}'\bfeps$ is asymptotically unit normal after standardizing
by the inverse square root of its asymptotic covariance matrix. The asymptotic covariance (dispersion) matrix of $\sqrt{n}(\hat{\bfbeta}_n - \bfbeta_o)$ is $\mathbf{D}_n$, which can be
consistently estimated by $\hat{\mathbf{D}}_n$. Note that this result allows the regressors
to be stochastic and imposes no restriction on the serial correlation or heteroskedasticity of $\eps_t$, except that needed to ensure that $(ii)$ holds. As we
shall see in the next chapter, only mild restrictions need to be imposed in
guaranteeing $(ii)$.

In special cases, it may be known that $\mathbf{V}_n$ has a special form. For example, when $\eps_t$ is an i.i.d.\ scalar with $E(\eps_t) = 0$, $E(\eps_t^2) = \sigma_o^2$, and $\mathbf{X}_t$ is
nonstochastic, then $\mathbf{V}_n = \sigma_o^2\mathbf{X}'\mathbf{X}/n$. Finding a consistent estimator for $\mathbf{V}_n$
then requires no more than finding a consistent estimator for $\sigma_o^2$.

In more general cases considered below it is often possible to write
\[
\mathbf{V}_n = E(\mathbf{X}'\bfeps\bfeps'\mathbf{X}/n) = E(\mathbf{X}'\bfOmega_n\mathbf{X}/n).
\]
Finding a consistent estimator for $\mathbf{V}_n$ in these cases is made easier by the
knowledge of the structure of $\bfOmega_n$. However, even when $\bfOmega_n$ is unknown, it
turns out that consistent estimators for $\mathbf{V}_n$ are generally available. The
conditions under which $\mathbf{V}_n$ can be consistently estimated are treated in
Chapter~6.

A result analogous to Theorem~4.25 is available for instrumental variables
estimators. Because the proof follows that of Theorem~4.25 very closely,
proof of the following result is left as an exercise.

\begin{exercises}
\item\label{ex:4.26} \textit{Prove the following result. Given}
\begin{enumerate}[label=(\roman*)]
\item $\mathbf{Y}_t = \mathbf{X}_t'\bfbeta_o + \eps_t$, \quad $t = 1, 2, \ldots$\,, $\bfbeta_o \in \mathbb{R}^k$;
\item $\mathbf{V}_n^{-1/2}n^{-1/2}\mathbf{Z}'\bfeps \overset{A}{\sim} N(\mathbf{0}, \mathbf{I})$, where $\mathbf{V}_n \equiv \mathrm{var}(n^{-1/2}\mathbf{Z}'\bfeps)$ is $O(1)$ and uniformly positive definite;
\item \begin{enumerate}[label=(\alph*)]
\item $\mathbf{Z}'\mathbf{X}/n - \mathbf{Q}_n \convp \mathbf{0}$, where $\mathbf{Q}_n \equiv E(\mathbf{Z}'\mathbf{X}/n)$ is $O(1)$ with uniformly full column rank;
\item There exists $\hat{\mathbf{P}}_n$ such that $\hat{\mathbf{P}}_n - \mathbf{P}_n \convp \mathbf{0}$ and $\mathbf{P}_n = O(1)$ and is symmetric and uniformly positive definite.
\end{enumerate}
\end{enumerate}
\emph{Then} $\mathbf{D}_n^{-1/2}\sqrt{n}(\tilde{\bfbeta}_n - \bfbeta_o) \overset{A}{\sim} N(\mathbf{0}, \mathbf{I})$, \emph{where}
\[
\mathbf{D}_n \equiv (\mathbf{Q}_n'\mathbf{P}_n\mathbf{Q}_n)^{-1}\mathbf{Q}_n'\mathbf{P}_n\mathbf{V}_n\mathbf{P}_n\mathbf{Q}_n(\mathbf{Q}_n'\mathbf{P}_n\mathbf{Q}_n)^{-1}
\]
\emph{and $\mathbf{D}_n^{-1}$ are $O(1)$.}

\emph{Suppose in addition that}
\begin{enumerate}[label=(\roman*)]
\setcounter{enumi}{3}
\item \emph{There exists a matrix $\hat{\mathbf{V}}_n$ positive semidefinite and symmetric such
that $\hat{\mathbf{V}}_n - \mathbf{V}_n \convp \mathbf{0}$.}

\emph{Then $\hat{\mathbf{D}}_n - \mathbf{D}_n \convp \mathbf{0}$, where}
\[
\hat{\mathbf{D}}_n = (\mathbf{X}'\mathbf{Z}\hat{\mathbf{P}}_n\mathbf{Z}'\mathbf{X}/n^2)^{-1}(\mathbf{X}'\mathbf{Z}/n)\hat{\mathbf{P}}_n\hat{\mathbf{V}}_n\hat{\mathbf{P}}_n(\mathbf{Z}'\mathbf{X}/n)(\mathbf{X}'\mathbf{Z}\hat{\mathbf{P}}_n\mathbf{Z}'\mathbf{X}/n^2)^{-1}.
\]
\end{enumerate}
\end{exercises}

\section{Hypothesis Testing}\label{sec:hypothesis-testing}

A direct and very important use of the asymptotic normality of a given
estimator is in hypothesis testing. Often, hypotheses of interest can be
expressed in terms of linear combinations of the parameters as
\[
\mathbf{R}\bfbeta_\bullet = \mathbf{r},
\]
where $\mathbf{R}$ is a given $q \times k$ matrix and $\mathbf{r}$ is a given $q \times 1$ vector that, through
$\mathbf{R}\bfbeta_o = \mathbf{r}$, specify the hypotheses of interest. For example, if the hypothesis
is that the elements of $\bfbeta_o$ sum to unity, $\mathbf{R} = [1, \ldots, 1]$ and $\mathbf{r} = 1$.

Several different approaches can be taken in computing a statistic to test
the null hypothesis $\mathbf{R}\bfbeta_o = \mathbf{r}$ versus the alternative $\mathbf{R}\bfbeta_o \neq \mathbf{r}$. The methods
that we consider here involve the use of Wald, Lagrange multiplier, and
quasi-likelihood ratio statistics.

Although the approaches to forming the test statistics differ, the way
that we determine their asymptotic distributions is the same. In each case
we exploit an underlying asymptotic normality property to obtain a statistic distributed asymptotically as chi-squared ($\chi^2$). To do this we use the
following results.

\begin{lemma}\label{lem:4.27}
Let $\mathbf{g} : \mathbb{R}^k \to \mathbb{R}^l$ be continuous on $\mathbb{R}^k$ and let $\mathbf{b}_n \convd \mathcal{Z}$, a $k \times 1$ random vector. Then $\mathbf{g}(\mathbf{b}_n) \to \mathbf{g}(\mathcal{Z})$.
\end{lemma}

\begin{proof}
See Rao (1973, p.~124).
\end{proof}

\begin{corollary}\label{cor:4.28}
Let $\mathbf{V}_n^{-1/2}\mathbf{b}_n \overset{A}{\sim} N(\mathbf{0}, \mathbf{I}_k)$. Then
\[
\mathbf{b}_n'\mathbf{V}_n^{-1}\mathbf{b}_n = \mathbf{b}_n'\mathbf{V}_n^{-1/2}\mathbf{V}_n^{-1/2}\mathbf{b}_n \overset{A}{\sim} \chi_k^2,
\]
where $\chi_k^2$ is a chi-squared random variable with $k$ degrees of freedom.
\end{corollary}

\begin{proof}
By hypothesis, $\mathbf{V}_n^{-1/2}\mathbf{b}_n \convd \mathcal{Z} \sim N(\mathbf{0}, \mathbf{I}_k)$. The function $g(\mathbf{z}) = \mathbf{z}'\mathbf{z}$ is continuous on $\mathbb{R}^k$. Hence,
\[
\mathbf{b}_n'\mathbf{V}_n^{-1}\mathbf{b}_n = g(\mathbf{V}^{-1/2}\mathbf{b}_n) \convd g(\mathcal{Z}) = \mathcal{Z}'\mathcal{Z} \sim \chi_k^2.
\]
\end{proof}

Typically, $\mathbf{V}_n$ will be unknown, but there will be a consistent estimator
$\hat{\mathbf{V}}_n$ such that $\hat{\mathbf{V}}_n - \mathbf{V}_n \convp \mathbf{0}$. To replace $\mathbf{V}_n$ in Corollary~4.28 with $\hat{\mathbf{V}}_n$,
we use the following result.

\begin{lemma}\label{lem:4.29}
Let $\mathbf{g} : \mathbb{R}^k \to \mathbb{R}^l$ be continuous on $\mathbb{R}^k$. If $\mathbf{a}_n - \mathbf{b}_n \convp \mathbf{0}$ and $\mathbf{b}_n \convd \mathcal{Z}$, then $\mathbf{g}(\mathbf{a}_n) - \mathbf{g}(\mathbf{b}_n) \convp \mathbf{0}$ and $\mathbf{g}(\mathbf{a}_n) \convd \mathbf{g}(\mathcal{Z})$.
\end{lemma}

\begin{proof}
Rao (1973, p.~124) proves that $\mathbf{g}(\mathbf{a}_n) - \mathbf{g}(\mathbf{b}_n) \convp \mathbf{0}$. That $\mathbf{g}(\mathbf{a}_n) \convd \mathbf{g}(\mathcal{Z})$ follows from Lemmas~4.7 and 4.27.
\end{proof}

Now we can prove the result that is the basis for finding the asymptotic
distribution of the Wald, Lagrange multiplier, and quasi-likelihood ratio
tests.

\begin{theorem}\label{thm:4.30}
Let $\mathbf{V}_n^{-1/2}\mathbf{b}_n \overset{A}{\sim} N(\mathbf{0}, \mathbf{I}_k)$, and suppose there exists $\hat{\mathbf{V}}_n$ positive semidefinite and symmetric such that $\hat{\mathbf{V}}_n - \mathbf{V}_n \convp \mathbf{0}$, where $\mathbf{V}_n$ is $O(1)$, and for all $n$ sufficiently large, $\det(\mathbf{V}_n) > \delta > 0$. Then $\mathbf{b}_n'\hat{\mathbf{V}}_n^{-1}\mathbf{b}_n \overset{A}{\sim} \chi_k^2$.
\end{theorem}

\begin{proof}
We apply Lemma~4.29. Consider $\hat{\mathbf{V}}_n^{-1/2}\mathbf{b}_n - \mathbf{V}_n^{-1/2}\mathbf{b}_n$, where $\hat{\mathbf{V}}_n^{-1/2}$
exists in probability for all $n$ sufficiently large. Now
\[
\hat{\mathbf{V}}_n^{-1/2}\mathbf{b}_n - \mathbf{V}_n^{-1/2}\mathbf{b}_n = (\hat{\mathbf{V}}_n^{-1/2}\mathbf{V}_n^{1/2} - \mathbf{I})\mathbf{V}_n^{-1/2}\mathbf{b}_n.
\]
By hypothesis $\mathbf{V}_n^{-1/2}\mathbf{b}_n \overset{A}{\sim} N(\mathbf{0}, \mathbf{I}_k)$ and $\hat{\mathbf{V}}_n^{-1/2}\mathbf{V}_n^{1/2} - \mathbf{I} \convp \mathbf{0}$ by Proposition~2.30. It follows from the product rule Lemma~4.6 that $\hat{\mathbf{V}}_n^{-1/2}\mathbf{b}_n -
\mathbf{V}_n^{-1/2}\mathbf{b}_n \convp \mathbf{0}$. Because $\mathbf{V}_n^{-1/2}\mathbf{b}_n \convd \mathcal{Z} \sim N(\mathbf{0}, \mathbf{I}_k)$, it follows from Lemma
4.29 that $\mathbf{b}_n'\hat{\mathbf{V}}_n^{-1}\mathbf{b}_n \overset{A}{\sim} \chi_k^2$.
\end{proof}

The Wald statistic allows the simplest analysis, although it may or may
not be the easiest statistic to compute in a given situation. The motivation
for the Wald statistic is that when the null hypothesis is correct, $\mathbf{R}\hat{\bfbeta}_n$
should be close to $\mathbf{R}\bfbeta_o = \mathbf{r}$, so a value of $\mathbf{R}\hat{\bfbeta}_n - \mathbf{r}$ far from zero is evidence
against the null hypothesis. To tell how far from zero $\mathbf{R}\hat{\bfbeta}_n - \mathbf{r}$ must be
before we reject the null hypothesis, we need to determine its asymptotic
distribution.

\begin{theorem}[Wald test]\label{thm:4.31}
Let the conditions of Theorem~4.25 hold and let $\mathrm{rank}(\mathbf{R}) = q < k$. Then under $H_0$: $\mathbf{R}\bfbeta_o = \mathbf{r}$,
\begin{enumerate}[label=(\roman*)]
\item $\boldsymbol{\Gamma}_n^{-1/2}\sqrt{n}(\mathbf{R}\hat{\bfbeta}_n - \mathbf{r}) \overset{A}{\sim} N(\mathbf{0}, \mathbf{I})$, where
\[
\boldsymbol{\Gamma}_n \equiv \mathbf{R}\mathbf{D}_n\mathbf{R}' = \mathbf{R}\mathbf{M}_n^{-1}\mathbf{V}_n\mathbf{M}_n^{-1}\mathbf{R}'.
\]

\item The Wald statistic $\mathcal{W}_n \equiv n(\mathbf{R}\hat{\bfbeta}_n - \mathbf{r})'\hat{\boldsymbol{\Gamma}}_n^{-1}(\mathbf{R}\hat{\bfbeta}_n - \mathbf{r}) \overset{A}{\sim} \chi_q^2$, where
\[
\hat{\boldsymbol{\Gamma}}_n \equiv \mathbf{R}\hat{\mathbf{D}}_n\mathbf{R}' = \mathbf{R}(\mathbf{X}'\mathbf{X}/n)^{-1}\hat{\mathbf{V}}_n(\mathbf{X}'\mathbf{X}/n)^{-1}\mathbf{R}'.
\]
\end{enumerate}
\end{theorem}

\begin{proof}
$(i)$ Under $H_0$, $\mathbf{R}\hat{\bfbeta}_n - \mathbf{r} = \mathbf{R}(\hat{\bfbeta}_n - \bfbeta_o)$, so
\[
\boldsymbol{\Gamma}_n^{-1/2}\sqrt{n}(\mathbf{R}\hat{\bfbeta}_n - \mathbf{r}) = \boldsymbol{\Gamma}_n^{-1/2}\mathbf{R}\mathbf{D}_n^{1/2}\mathbf{D}_n^{-1/2}\sqrt{n}(\hat{\bfbeta}_n - \bfbeta_o).
\]
It follows from Corollary~4.24 that $\boldsymbol{\Gamma}_n^{-1/2}\sqrt{n}(\mathbf{R}\hat{\bfbeta}_n - \mathbf{r}) \overset{A}{\sim} N(\mathbf{0}, \mathbf{I})$.

$(ii)$ Because $\hat{\mathbf{D}}_n - \mathbf{D}_n \convp \mathbf{0}$ from Theorem~4.25, it follows from Proposition~2.30 that $\hat{\boldsymbol{\Gamma}}_n - \boldsymbol{\Gamma}_n \convp \mathbf{0}$. Given the result in $(i)$, $(ii)$ follows from
Theorem~4.30.
\end{proof}

This version of the Wald statistic is useful regardless of the presence of
heteroskedasticity or serial correlation in the error terms because a consistent estimator ($\hat{\mathbf{V}}_n$) for $\mathbf{V}_n$ is used in computing $\hat{\boldsymbol{\Gamma}}_n$. In the special case
when $\mathbf{V}_n$ can be consistently estimated by $\hat{\sigma}_n^2(\mathbf{X}'\mathbf{X}/n)$, the Wald test has
the form
\[
\mathcal{W}_n = n(\mathbf{R}\hat{\bfbeta}_n - \mathbf{r})'[\mathbf{R}(\mathbf{X}'\mathbf{X}/n)^{-1}\mathbf{R}']^{-1}(\mathbf{R}\hat{\bfbeta}_n - \mathbf{r})/\hat{\sigma}_n^2,
\]
which is simply $q$ times the standard $F$-statistic for testing the hypothesis
$\mathbf{R}\bfbeta_o = \mathbf{r}$. The validity of the asymptotic $\chi_q^2$ distribution for this statistic
depends crucially on the consistency of $\hat{\mathbf{V}}_n = \hat{\sigma}_n^2(\mathbf{X}'\mathbf{X}/n)$ for $\mathbf{V}_n$; if this
$\hat{\mathbf{V}}_n$ is not consistent for $\mathbf{V}_n$, the asymptotic distribution of this form for
$\mathcal{W}_n$ is not $\chi_q^2$ in general.

The Wald statistic is most convenient in situations in which the restrictions $\mathbf{R}\bfbeta_o = \mathbf{r}$ are not easy to impose in estimating $\bfbeta_o$. When these restrictions are easily imposed (say, $\mathbf{R}\bfbeta_o = \mathbf{r}$ specifies that the last element
of $\bfbeta_o$ is zero), the Lagrange multiplier statistic is more easily computed.

The motivation for the Lagrange multiplier statistic is that a constrained
least squares estimator can be obtained by solving the problem
\[
\min_{\bfbeta} (\mathbf{Y} - \mathbf{X}\bfbeta)'(\mathbf{Y} - \mathbf{X}\bfbeta)/n, \qquad \text{s.t.}\ \mathbf{R}\bfbeta = \mathbf{r},
\]
which is equivalent to finding the saddle point of the Lagrangian
\[
\mathcal{L} = (\mathbf{Y} - \mathbf{X}\bfbeta)'(\mathbf{Y} - \mathbf{X}\bfbeta)/n + (\mathbf{R}\bfbeta - \mathbf{r})'\boldsymbol{\lambda}.
\]
The Lagrange multipliers $\boldsymbol{\lambda}$ can be thought of as giving the shadow
price of the constraint and should therefore be small when the constraint is
valid and large otherwise. (See Engle, 1981, for a general discussion.) The
Lagrange multiplier test can be thought of as testing the hypothesis that
$\boldsymbol{\lambda} = \mathbf{0}$.

The first order conditions are
\begin{eqnarray*}
\partial \mathcal{L}/\partial \bfbeta &=& 2(\mathbf{X}'\mathbf{X}/n)\bfbeta - 2\mathbf{X}'\mathbf{Y}/n + \mathbf{R}'\boldsymbol{\lambda} = \mathbf{0}\\
\partial \mathcal{L}/\partial \boldsymbol{\lambda} &=& \mathbf{R}\bfbeta - \mathbf{r} = \mathbf{0}.
\end{eqnarray*}

To solve for the estimate of the Lagrange multiplier, premultiply the first
equation by $\mathbf{R}(\mathbf{X}'\mathbf{X}/n)^{-1}$ and set $\mathbf{R}\bfbeta = \mathbf{r}$. This yields
\begin{eqnarray*}
\ddot{\boldsymbol{\lambda}}_n &=& 2(\mathbf{R}(\mathbf{X}'\mathbf{X}/n)^{-1}\mathbf{R}')^{-1}(\mathbf{R}\hat{\bfbeta}_n - \mathbf{r})\\
\ddot{\bfbeta}_n &=& \hat{\bfbeta}_n - (\mathbf{X}'\mathbf{X}/n)^{-1}\mathbf{R}'\ddot{\boldsymbol{\lambda}}_n/2,
\end{eqnarray*}
where $\ddot{\bfbeta}_n$ is the constrained least squares estimator (which automatically
satisfies $\mathbf{R}\ddot{\bfbeta}_n = \mathbf{r}$). In this form, $\ddot{\boldsymbol{\lambda}}_n$ is simply a nonsingular transformation
of $\mathbf{R}\hat{\bfbeta}_n - \mathbf{r}$. This allows the following result to be proved very simply.

\begin{theorem}[Lagrange multiplier test]\label{thm:4.32}
Let the conditions of Theorem~4.25 hold and let $\mathrm{rank}(\mathbf{R}) = q < k$. Then under $H_0$: $\mathbf{R}\bfbeta_o = \mathbf{r}$,
\begin{enumerate}[label=(\roman*)]
\item $\boldsymbol{\Lambda}_n^{-1/2}\sqrt{n}\ddot{\boldsymbol{\lambda}}_n \overset{A}{\sim} N(\mathbf{0}, \mathbf{I})$, where
\[
\boldsymbol{\Lambda}_n \equiv 4(\mathbf{R}\mathbf{M}_n^{-1}\mathbf{R}')^{-1}\boldsymbol{\Gamma}_n(\mathbf{R}\mathbf{M}_n^{-1}\mathbf{R}')^{-1}
\]
and $\boldsymbol{\Gamma}_n$ is as defined in Theorem~4.31.

\item The Lagrange multiplier statistic $\mathcal{LM}_n \equiv n\ddot{\boldsymbol{\lambda}}_n'\hat{\boldsymbol{\Lambda}}_n^{-1}\ddot{\boldsymbol{\lambda}}_n \overset{A}{\sim} \chi_q^2$, where
\begin{align*}
\hat{\boldsymbol{\Lambda}}_n &\equiv 4(\mathbf{R}(\mathbf{X}'\mathbf{X}/n)^{-1}\mathbf{R}')^{-1}\mathbf{R}(\mathbf{X}'\mathbf{X}/n)^{-1}\ddot{\mathbf{V}}_n(\mathbf{X}'\mathbf{X}/n)^{-1}\mathbf{R}'\\
&\quad\times(\mathbf{R}(\mathbf{X}'\mathbf{X}/n)^{-1}\mathbf{R}')^{-1}
\end{align*}
and $\ddot{\mathbf{V}}_n$ is computed from the constrained regression such that $\ddot{\mathbf{V}}_n - \mathbf{V}_n \convp \mathbf{0}$ under $H_0$.
\end{enumerate}
\end{theorem}

\begin{proof}
$(i)$ Consider the difference
\begin{align*}
&\boldsymbol{\Lambda}_n^{-1/2}\sqrt{n}\ddot{\boldsymbol{\lambda}}_n - 2\boldsymbol{\Lambda}_n^{-1/2}(\mathbf{R}\mathbf{M}_n^{-1}\mathbf{R}')^{-1}\sqrt{n}(\mathbf{R}\hat{\bfbeta}_n - \mathbf{r})\\
&= 2\boldsymbol{\Lambda}_n^{-1/2}[(\mathbf{R}(\mathbf{X}'\mathbf{X}/n)^{-1}\mathbf{R}')^{-1} - (\mathbf{R}\mathbf{M}_n^{-1}\mathbf{R}')^{-1}]\boldsymbol{\Gamma}_n^{1/2}\boldsymbol{\Gamma}_n^{-1/2}\sqrt{n}(\mathbf{R}\hat{\bfbeta}_n - \mathbf{r}).
\end{align*}
From Theorem~4.31, $\boldsymbol{\Gamma}_n^{-1/2}\sqrt{n}(\mathbf{R}\hat{\bfbeta}_n - \mathbf{r}) \overset{A}{\sim} N(\mathbf{0}, \mathbf{I})$. Because $(\mathbf{X}'\mathbf{X}/n) -
\mathbf{M}_n \convp \mathbf{0}$, it follows from Proposition~2.30 and the fact that $\boldsymbol{\Lambda}_n^{-1/2}$ and
$\boldsymbol{\Gamma}_n^{1/2}$ are $O(1)$ that
\[
\boldsymbol{\Lambda}_n^{-1/2}\sqrt{n}\ddot{\boldsymbol{\lambda}}_n - 2\boldsymbol{\Lambda}_n^{-1/2}(\mathbf{R}\mathbf{M}_n^{-1}\mathbf{R}')^{-1}\sqrt{n}(\mathbf{R}\hat{\bfbeta}_n - \mathbf{r}) \convp \mathbf{0}.
\]
Hence by the product rule Lemma~4.6,
$\boldsymbol{\Lambda}_n^{-1/2}\sqrt{n}\ddot{\boldsymbol{\lambda}}_n - 2\boldsymbol{\Lambda}_n^{-1/2}(\mathbf{R}\mathbf{M}_n^{-1}\mathbf{R}')^{-1}\sqrt{n}(\mathbf{R}\hat{\bfbeta}_n - \mathbf{r}) \convp \mathbf{0}$.
It follows from Lemma~4.7 that $\boldsymbol{\Lambda}_n^{-1/2}\sqrt{n}\ddot{\boldsymbol{\lambda}}_n$ has the same asymptotic distribution as $2\boldsymbol{\Lambda}_n^{-1/2}(\mathbf{R}\mathbf{M}_n^{-1}\mathbf{R}')^{-1}\sqrt{n}(\mathbf{R}\hat{\bfbeta}_n - \mathbf{r})$. It follows immediately from
Corollary~4.24 that $2\boldsymbol{\Lambda}_n^{-1/2}(\mathbf{R}\mathbf{M}_n^{-1}\mathbf{R}')^{-1}\sqrt{n}(\mathbf{R}\bfbeta_n - \mathbf{r}) \overset{A}{\sim} N(\mathbf{0}, \mathbf{I})$; hence
$\boldsymbol{\Lambda}_n^{-1/2}\sqrt{n}\ddot{\boldsymbol{\lambda}}_n \overset{A}{\sim} N(\mathbf{0}, \mathbf{I})$.

$(ii)$ Because $\ddot{\mathbf{V}}_n - \mathbf{V}_n \convp \mathbf{0}$ by hypothesis and $(\mathbf{X}'\mathbf{X}/n) - \mathbf{M}_n \convp \mathbf{0}$,
$\hat{\boldsymbol{\Lambda}}_n - \boldsymbol{\Lambda}_n \convp \mathbf{0}$ by Proposition~2.30, given the result in $(i)$, $(ii)$ follows from
Theorem~4.30.
\end{proof}

Note that the Wald and Lagrange multiplier statistics would be identical
if $\hat{\mathbf{V}}_n$ were used in place of $\ddot{\mathbf{V}}_n$. This suggests that the two statistics should
be asymptotically equivalent.

\begin{exercises}
\item\label{ex:4.33} \textit{Prove that under the conditions of Theorems~4.31 and 4.32, $\mathcal{W}_n - \mathcal{LM}_n \convp 0$.}
\end{exercises}

Although the fact that $\ddot{\boldsymbol{\lambda}}_n$ is a linear combination of $\mathbf{R}\hat{\bfbeta}_n - \mathbf{r}$ simplifies the
proof of Theorem~4.32, the whole point of using the Lagrange multiplier
statistic is to avoid computing $\hat{\bfbeta}_n$ and to compute only the simpler $\ddot{\bfbeta}_n$.
Computation of $\ddot{\bfbeta}_n$ is particularly simple when the data are generated as
$\mathbf{Y} = \mathbf{X}_1\bfbeta_1 + \mathbf{X}_2\bfbeta_2 + \bfeps$ and $H_0$ specifies that $\bfbeta_2$ (a $q \times 1$ vector) is zero.
Then
\[
\mathbf{R} = [\underset{(q\times(k-q))}{\mathbf{0}} \vdots \underset{(q\times q)}{\mathbf{I}}], \qquad \mathbf{r} = \underset{(q\times 1)}{\mathbf{0}},
\]
and $\ddot{\bfbeta}_n' = (\ddot{\bfbeta}_{1n}', \mathbf{0})$ where $\ddot{\bfbeta}_{1n} = (\mathbf{X}_1'\mathbf{X}_1)^{-1}\mathbf{X}_1'\mathbf{Y}$.

\begin{exercises}
\item\label{ex:4.34} \textit{Define $\ddot{\bfeps} = \mathbf{Y} - \mathbf{X}_1\ddot{\bfbeta}_{1n}$. Show that under $H_0$: $\bfbeta_2 = \mathbf{0}$,}
\[
\ddot{\boldsymbol{\lambda}}_n = 2\mathbf{X}_2'(\mathbf{I} - \mathbf{X}_1(\mathbf{X}_1'\mathbf{X}_1)^{-1}\mathbf{X}_1')\bfeps/n = 2\mathbf{X}_2'\ddot{\bfeps}/n.
\]
\emph{(Hint: $\mathbf{R}\hat{\bfbeta}_n - \mathbf{r} = \mathbf{R}(\mathbf{X}'\mathbf{X}/n)^{-1}\mathbf{X}'(\mathbf{Y} - \mathbf{X}\ddot{\bfbeta}_n)/n$.)}

\item\label{ex:4.35} \textit{If $\ddot{\sigma}_n^2(\mathbf{X}'\mathbf{X}/n) - \mathbf{V}_n \convp \mathbf{0}$ and $\bfbeta_2 = \mathbf{0}$, show that $\mathcal{LM}_n = n\ddot{\bfeps}'\mathbf{X}(\mathbf{X}'\mathbf{X})^{-1}\mathbf{X}'\ddot{\bfeps}/(\ddot{\bfeps}'\ddot{\bfeps})$, which is $n$ times the simple $R^2$ of the regression of $\ddot{\bfeps}$ on $\mathbf{X}$.}
\end{exercises}

When $\mathbf{V}_n$ can be consistently estimated by $\ddot{\mathbf{V}}_n = \ddot{\sigma}_n^2(\mathbf{X}'\mathbf{X}/n)$, where $\ddot{\sigma}_n^2 = \ddot{\bfeps}'\ddot{\bfeps}/n$, the $\mathcal{LM}_n$ statistic simplifies even further.

The result of this exercise implies a very simple procedure for testing
$\bfbeta_2 = \mathbf{0}$ when $\mathbf{V}_n = \sigma_o^2\mathbf{M}_n$. First, regress $\mathbf{Y}$ on $\mathbf{X}_1$ and form the constrained
residuals $\ddot{\bfeps}$. Then regress $\ddot{\bfeps}$ on $\mathbf{X}$. The product of the sample size $n$ and the
simple $R^2$ (i.e., without adjustment for the presence of a constant in the
regression) from this regression is the $\mathcal{LM}_n$ test statistic, which has the
$\chi_q^2$ distribution asymptotically. As Engle (1981) showed, many interesting
diagnostic statistics can be computed in this way.

When the errors $\eps_t$ are scalar i.i.d.\ $N(0, \sigma_o^2)$ random variables, the OLS
estimator is also the maximum likelihood estimator (MLE) because $\hat{\bfbeta}_n$
solves the problem
\[
\max \mathcal{L}(\bfbeta, \sigma; \mathbf{Y}) = \exp[-n\log\sqrt{2\pi} - n\log\sigma - \frac{1}{2}\sum_{t=1}^{n}(Y_t - \mathbf{X}_t'\bfbeta)^2/\sigma^2],
\]
where $\mathcal{L}(\bfbeta, \sigma; \mathbf{Y})$ is the sample likelihood based on the normality assumption. When $\eps_t$ is not i.i.d.\ $N(0, \sigma_o^2)$, $\hat{\bfbeta}_n$ is said to be a quasi-maximum
likelihood estimator (QMLE).

When $\hat{\bfbeta}_n$ is the MLE, hypothesis tests can be based on the log-likelihood
ratio
\[
\mathcal{LR}_n = \log\left[\frac{\mathcal{L}(\ddot{\bfbeta}_n, \ddot{\sigma}_n; \mathbf{Y})}{\mathcal{L}(\hat{\bfbeta}_n, \hat{\sigma}_n; \mathbf{Y})}\right],
\]
where $\hat{\sigma}_n^2 = n^{-1}\sum_{t=1}^{n}(Y_t - \mathbf{X}_t'\hat{\bfbeta}_n)^2$ as before and $\ddot{\bfbeta}_n$, $\ddot{\sigma}_n$ solve
\[
\max \mathcal{L}(\bfbeta, \sigma; \mathbf{Y}), \qquad \text{s.t.}\ \mathbf{R}\bfbeta = \mathbf{r}.
\]
It is easy to show that $\ddot{\bfbeta}_n$ is the constrained OLS estimator and $\ddot{\sigma}_n^2 = \ddot{\bfeps}'\ddot{\bfeps}/n$.
The likelihood ratio is nonnegative and always less than or equal
to 1. Simple algebra yields
\[
\mathcal{LR}_n = (n/2)\log(\hat{\sigma}_n^2/\ddot{\sigma}_n^2).
\]
Because $\ddot{\sigma}_n^2 = \hat{\sigma}_n^2 + (\hat{\bfbeta}_n - \ddot{\bfbeta}_n)'(\mathbf{X}'\mathbf{X}/n)(\hat{\bfbeta}_n - \ddot{\bfbeta}_n)$ (verify this),
\[
\mathcal{LR}_n = -(n/2)\log[1 + (\hat{\bfbeta}_n - \ddot{\bfbeta}_n)'(\mathbf{X}'\mathbf{X}/n)(\hat{\bfbeta}_n - \ddot{\bfbeta}_n)/\hat{\sigma}_n^2].
\]

To find the asymptotic distribution of this statistic, we make use of the
mean value theorem of calculus.

\begin{theorem}[Mean value theorem]\label{thm:4.36}
Let $s : \mathbb{R}^k \to \mathbb{R}$ be defined on an open convex set $\Theta \subset \mathbb{R}^k$ such that $s$ is continuously differentiable on $\Theta$ with $k \times 1$ gradient $\nabla s$. Then for any points $\boldsymbol{\theta}$ and $\boldsymbol{\theta}_o \in \Theta$ there exists $\bar{\boldsymbol{\theta}}$ on the segment connecting $\boldsymbol{\theta}$ and $\boldsymbol{\theta}_o$ such that $s(\boldsymbol{\theta}) = s(\boldsymbol{\theta}_o) + \nabla s(\bar{\boldsymbol{\theta}})'(\boldsymbol{\theta} - \boldsymbol{\theta}_o)$.
\end{theorem}

\begin{proof}
See Bartle (1976, p.~365).
\end{proof}

For the present application, we choose $s(\theta) = \log(1+\theta)$. If we also choose
$\theta_o = 0$, we have $s(\theta) = \log(1)+(1/(1+\bar{\theta}))\theta = \theta/(1+\bar{\theta})$, where $\bar{\theta}$ lies between
$\theta$ and zero. Let $\theta_n = (\hat{\bfbeta}_n - \ddot{\bfbeta}_n)'(\mathbf{X}'\mathbf{X}/n)(\hat{\bfbeta}_n - \ddot{\bfbeta}_n)/\ddot{\sigma}_n^2$ so that under $H_0$,
$|\bar{\theta}_n| < |\theta_n| \convp 0$; hence $\bar{\theta}_n \convp 0$. Applying the mean value theorem now
gives
\[
\mathcal{LR}_n = -(n/2)(1+\bar{\theta}_n)^{-1}(\hat{\bfbeta}_n - \ddot{\bfbeta}_n)'(\mathbf{X}'\mathbf{X}/n)(\hat{\bfbeta}_n - \ddot{\bfbeta}_n)/\hat{\sigma}_n^2.
\]
Because $(1+\bar{\theta}_n)^{-1} \convp 1$, it follows from Lemma~4.6 that
\[
-2\mathcal{LR}_n - n(\hat{\bfbeta}_n - \ddot{\bfbeta}_n)'(\mathbf{X}'\mathbf{X}/n)(\hat{\bfbeta}_n - \ddot{\bfbeta}_n)/\hat{\sigma}_n^2 \convp 0,
\]
provided the second term has a limiting distribution. Now
\[
\hat{\bfbeta}_n - \ddot{\bfbeta}_n = (\mathbf{X}'\mathbf{X}/n)^{-1}\mathbf{R}'(\mathbf{R}(\mathbf{X}'\mathbf{X}/n)^{-1}\mathbf{R}')^{-1}(\mathbf{R}\hat{\bfbeta}_n - \mathbf{r}).
\]
Thus
\[
-2\mathcal{LR}_n - n(\mathbf{R}\hat{\bfbeta}_n - \mathbf{r})'[\mathbf{R}(\mathbf{X}'\mathbf{X}/n)^{-1}\mathbf{R}']^{-1}(\mathbf{R}\hat{\bfbeta}_n - \mathbf{r})/\hat{\sigma}_n^2 \convp 0.
\]
This second term is the Wald statistic formed with $\hat{\mathbf{V}}_n = \hat{\sigma}_n^2(\mathbf{X}'\mathbf{X}/n)$, so
$-2\mathcal{LR}_n$ is asymptotically equivalent to the Wald statistic and has the $\chi_q^2$
distribution asymptotically, provided $\hat{\sigma}_n^2(\mathbf{X}'\mathbf{X}/n)$ is a consistent estimator
for $\mathbf{V}_n$. If this is not true, then $-2\mathcal{LR}_n$ does not in general have the $\chi_q^2$
distribution asymptotically. It does have a limiting distribution, but not a
simple one that has been tabulated or is easily computable (see White, 1994,
Ch.~6, for further details). Note that it is not violation of the normality
assumption per se, but the failure of $\mathbf{V}_n$ to equal $\sigma_o^2\mathbf{M}_n$ that results in
$-2\mathcal{LR}_n$ not having the $\chi_q^2$ distribution asymptotically.

The formal statement of the result for the $\mathcal{LR}_n$ statistic is the following.

\begin{theorem}[Likelihood ratio test]\label{thm:4.37}
Let the conditions of Theorem~4.25 hold, let $\mathrm{rank}(\mathbf{R}) = q < k$, and let $\hat{\sigma}_n^2(\mathbf{X}'\mathbf{X}/n) - \mathbf{V}_n \convp \mathbf{0}$. Then under $H_0$: $\mathbf{R}\bfbeta_o = \mathbf{r}$, $-2\mathcal{LR}_n \overset{A}{\sim} \chi_q^2$.
\end{theorem}

\begin{proof}
Set $\hat{\mathbf{V}}_n$ in Theorem~4.31 to $\hat{\mathbf{V}}_n = \hat{\sigma}_n^2(\mathbf{X}'\mathbf{X}/n)$. Then from the
argument preceding the theorem above $-2\mathcal{LR}_n - \mathcal{W}_n \convp 0$. Because $\mathcal{W}_n \overset{A}{\sim}
\chi_q^2$, it follows from Lemma~4.7 that $-2\mathcal{LR}_n \overset{A}{\sim} \chi_q^2$.
\end{proof}

The mean value theorem just introduced provides a convenient way to
find the asymptotic distribution of statistics used to test nonlinear hypotheses. In general, nonlinear hypotheses can be conveniently represented
as
\[
H_0:\ \mathbf{s}(\bfbeta_o) = \mathbf{0},
\]
where $\mathbf{s}: \mathbb{R}^k \to \mathbb{R}^q$ is a continuously differentiable function of $\bfbeta$.

\begin{example}\label{ex:4.38}
Suppose $\mathbf{Y} = \mathbf{X}_1\beta_1 + \mathbf{X}_2\beta_2 + \mathbf{X}_3\beta_3 + \bfeps$, where $\mathbf{X}_1$, $\mathbf{X}_2$, and $\mathbf{X}_3$ are $n \times 1$ and $\beta_1$, $\beta_2$, and $\beta_3$ are scalars. Further, suppose we hypothesize that $\beta_3 = \beta_1\beta_2$. Then $s(\bfbeta_o) = \beta_3 - \beta_1\beta_2 = 0$ expresses the null hypothesis.
\end{example}

Just as with linear restrictions, we can construct a Wald test based on the
asymptotic distribution of $\mathbf{s}(\hat{\bfbeta}_n)$; we can construct a Lagrange multiplier
test based on the Lagrange multipliers derived from minimizing the least
squares (or other estimation) objective function subject to the constraint;
or we can form a log-likelihood ratio.

To illustrate the approach, consider the Wald test based on $\mathbf{s}(\hat{\bfbeta}_n)$. As
before, a value of $\mathbf{s}(\hat{\bfbeta}_n)$ far from zero is evidence against $H_0$. To tell how far
$\mathbf{s}(\hat{\bfbeta}_n)$ must be from zero to reject $H_0$, we need to determine its asymptotic
distribution. This is provided by the next result.

\begin{theorem}[Wald test]\label{thm:4.39}
Let the conditions of Theorem~4.25 hold and let $\mathrm{rank}\,\nabla \mathbf{s}(\bfbeta_o) = q < k$, where $\nabla \mathbf{s}$ is the $k \times q$ gradient matrix of $\mathbf{s}$. Then under $H_0$: $\mathbf{s}(\bfbeta_o) = \mathbf{0}$,
\begin{enumerate}[label=(\roman*)]
\item $\boldsymbol{\Gamma}_n^{-1/2}\sqrt{n}\mathbf{s}(\hat{\bfbeta}_n) \overset{A}{\sim} N(\mathbf{0}, \mathbf{I})$, where
\[
\boldsymbol{\Gamma}_n \equiv \nabla \mathbf{s}(\bfbeta_o)'\mathbf{D}_n\nabla \mathbf{s}(\bfbeta_o).
\]

\item The Wald statistic $\mathcal{W}_n \equiv n\mathbf{s}(\hat{\bfbeta}_n)'\hat{\boldsymbol{\Gamma}}_n^{-1}\mathbf{s}(\hat{\bfbeta}_n) \overset{A}{\sim} \chi_q^2$, where
\begin{align*}
\hat{\boldsymbol{\Gamma}}_n &= \nabla \mathbf{s}(\hat{\bfbeta}_n)'\hat{\mathbf{D}}_n\nabla \mathbf{s}(\hat{\bfbeta}_n) \\
&= \nabla \mathbf{s}(\hat{\bfbeta}_n)'(\mathbf{X}'\mathbf{X}/n)^{-1}\hat{\mathbf{V}}_n(\mathbf{X}'\mathbf{X}/n)^{-1}\nabla \mathbf{s}(\hat{\bfbeta}_n).
\end{align*}
\end{enumerate}
\end{theorem}

\begin{proof}
$(i)$ Because $\mathbf{s}(\bfbeta)$ is a vector function, we apply the mean value
theorem to each element $s_i(\bfbeta)$, $i = 1, \ldots, q$, to get
\[
s_i(\hat{\bfbeta}_n) = s_i(\bfbeta_o) + \nabla s_i(\bar{\bfbeta}_n^{(i)})'(\hat{\bfbeta}_n - \bfbeta_o),
\]
where $\bar{\bfbeta}_n^{(i)}$ is a $k \times 1$ vector lying on the segment connecting $\hat{\bfbeta}_n$ and $\bfbeta_o$.
The superscript $(i)$ reflects the fact that the mean value may be different
for each element $s_i(\bfbeta)$ of $\mathbf{s}(\bfbeta)$.

Under $H_0$, $s_i(\bfbeta_o) = 0$, $i = 1, \ldots, q$, so
\[
\sqrt{n}s_i(\hat{\bfbeta}_n) = \nabla s_i(\bar{\bfbeta}_n^{(i)})'\sqrt{n}(\hat{\bfbeta}_n - \bfbeta_o).
\]
This suggests considering the difference
\begin{align*}
&\sqrt{n}s_i(\hat{\bfbeta}_n) - \nabla s_i(\bfbeta_o)'\sqrt{n}(\hat{\bfbeta}_n - \bfbeta_o) \\
&= (\nabla s_i(\bar{\bfbeta}_n^{(i)}) - \nabla s_i(\bfbeta_o))'\sqrt{n}(\hat{\bfbeta}_n - \bfbeta_o) \\
&= (\nabla s_i(\bar{\bfbeta}_n^{(i)}) - \nabla s_i(\bfbeta_o))'\mathbf{D}_n^{1/2}\mathbf{D}_n^{-1/2}\sqrt{n}(\hat{\bfbeta}_n - \bfbeta_o).
\end{align*}
By Theorem~4.25, $\mathbf{D}_n^{-1/2}\sqrt{n}(\hat{\bfbeta}_n - \bfbeta_o) \overset{A}{\sim} N(\mathbf{0}, \mathbf{I})$. Because $\hat{\bfbeta}_n \convp \bfbeta_o$, it
follows that $\bar{\bfbeta}_n^{(i)} \convp \bfbeta_o$ so $\nabla s_i(\bar{\bfbeta}_n^{(i)}) - \nabla s_i(\bfbeta_o) \convp \mathbf{0}$ by Proposition~2.27.
Because $\mathbf{D}_n^{1/2}$ is $O(1)$, we have $(\nabla s_i(\bar{\bfbeta}_n^{(i)}) - \nabla s_i(\bfbeta_o))'\mathbf{D}_n^{1/2} \convp \mathbf{0}$. It follows
from Lemma~4.6 that
\[
\sqrt{n}s_i(\hat{\bfbeta}_n) - \nabla s_i(\bfbeta_o)'\sqrt{n}(\hat{\bfbeta}_n - \bfbeta_o) \convp 0, \quad i = 1, \ldots, q.
\]
In vector form this becomes
\[
\sqrt{n}\mathbf{s}(\hat{\bfbeta}_n) - \nabla \mathbf{s}(\bfbeta_o)'\sqrt{n}(\hat{\bfbeta}_n - \bfbeta_o) \convp \mathbf{0},
\]
and because $\boldsymbol{\Gamma}_n^{-1/2}$ is $O(1)$,
\[
\boldsymbol{\Gamma}_n^{-1/2}\sqrt{n}\mathbf{s}(\hat{\bfbeta}_n) - \boldsymbol{\Gamma}_n^{-1/2}\nabla \mathbf{s}(\bfbeta_o)'\sqrt{n}(\hat{\bfbeta}_n - \bfbeta_o) \convp \mathbf{0}.
\]
Corollary~4.24 immediately yields
\[
\boldsymbol{\Gamma}_n^{-1/2}\nabla \mathbf{s}(\bfbeta_o)'\sqrt{n}(\hat{\bfbeta}_n - \bfbeta_o) \overset{A}{\sim} N(\mathbf{0}, \mathbf{I}),
\]
so by Lemma~4.7, $\boldsymbol{\Gamma}_n^{-1/2}\sqrt{n}\mathbf{s}(\hat{\bfbeta}_n) \overset{A}{\sim} N(\mathbf{0}, \mathbf{I})$.

$(ii)$ Because $\hat{\mathbf{D}}_n - \mathbf{D}_n \convp \mathbf{0}$ from Theorem~4.25 and $\nabla \mathbf{s}(\hat{\bfbeta}_n) - \nabla \mathbf{s}(\bfbeta_o) \convp
\mathbf{0}$ by Proposition~2.27, $\hat{\boldsymbol{\Gamma}}_n - \boldsymbol{\Gamma}_n \convp \mathbf{0}$ by Proposition~2.30. Given the result
in $(i)$, $(ii)$ follows from Theorem~4.30.
\end{proof}

Note the similarity of this result to Theorem~4.31 which gives the Wald
test for the linear hypothesis $\mathbf{R}\bfbeta_o = \mathbf{r}$. In the present context, $\mathbf{s}(\bfbeta_o)$ plays
the role of $\mathbf{R}\bfbeta_o - \mathbf{r}$, whereas $\nabla \mathbf{s}(\bfbeta_o)'$ plays the role of $\mathbf{R}$ in computing the
covariance matrix.

\begin{exercises}
\item\label{ex:4.40} \textit{Write down the Wald statistic for testing the hypothesis of
Example~4.38.}

\item\label{ex:4.41} \textit{Give the Lagrange multiplier statistic for testing the hypothesis $H_0$: $\mathbf{s}(\bfbeta_o) = \mathbf{0}$ versus $H_1$: $\mathbf{s}(\bfbeta_o) \neq \mathbf{0}$, and derive its limiting distribution under the conditions of Theorem~4.25.}

\item\label{ex:4.42} \textit{Give the Wald and Lagrange multiplier statistics for testing
the hypotheses $\mathbf{R}\bfbeta_o = \mathbf{r}$ and $\mathbf{s}(\bfbeta_o) = \mathbf{0}$ on the basis of the IV estimator $\tilde{\bfbeta}_n$ and derive their limiting distributions under the conditions of Exercise~4.26.}
\end{exercises}

\section{Asymptotic Efficiency}\label{sec:asymptotic-efficiency}

Given a class of estimators (e.g., the class of instrumental variables estimators), it is desirable to choose that member of the class that has the
smallest asymptotic covariance matrix (assuming that this member exists
and can be computed). The reason for this is that such estimators are obviously more precise, and in general allow construction of more powerful
test statistics. In what follows, we shall abuse notation slightly and write
$\mathrm{avar}(\tilde{\bfbeta}_n)$ instead of $\mathrm{avar}\left(\sqrt{n}(\tilde{\bfbeta}_n - \bfbeta_o)\right)$. A definition of asymptotic efficiency adequate for our purposes here is the following.

\begin{definition}\label{def:4.43}
Let $\mathcal{P}$ be a set of data generating processes such that each data generating process $P^o \in \mathcal{P}$ has a corresponding coefficient vector $\bfbeta^o \in \mathbb{R}^k$ and let $\mathcal{E}$ be a class of estimators $\{\bar{\bfbeta}_n\}$ such that for each $P^o \in \mathcal{P}$, $\mathbf{D}_n^{o\,-1/2}\sqrt{n}\left(\bar{\bfbeta}_n - \bfbeta^o\right) \xrightarrow{d^o} N(\mathbf{0}, \mathbf{I})$ for $\{\mathrm{avar}^o(\bar{\bfbeta}_n) \equiv \mathbf{D}_n^o\}$ nonstochastic, $O(1)$ and uniformly nonsingular. Then $\{\bfbeta_n^*\} \in \mathcal{E}$ is asymptotically efficient relative to $\{\bar{\bfbeta}_n\} \in \mathcal{E}$ for $\mathcal{P}$ if for every $P^o \in \mathcal{P}$ the matrix $\mathrm{avar}^o(\bar{\bfbeta}_n) - \mathrm{avar}^o(\bfbeta_n^*)$ is positive semidefinite for all $n$ sufficiently large. The estimator $\{\bfbeta_n^*\} \in \mathcal{E}$ is asymptotically efficient in $\mathcal{E}$ for $\mathcal{P}$ if it is asymptotically efficient relative to every other member of its class, $\mathcal{E}$, for $\mathcal{P}$.
\end{definition}

We write $\bfbeta^o$ instead of $\bfbeta_o$ to emphasize that $\bfbeta^o$ corresponds to (at least)
one possible data generating process ($P^o$) among many ($\mathcal{P}$).

The class of estimators we consider is the class $\mathcal{E}$ of instrumental variables
estimators
\[
\tilde{\bfbeta}_n = (\mathbf{X}'\mathbf{Z}\hat{\mathbf{P}}_n\mathbf{Z}'\mathbf{X})^{-1}\mathbf{X}'\mathbf{Z}\hat{\mathbf{P}}_n\mathbf{Z}'\mathbf{Y},
\]
where different members of the class are defined by the different possible
choices for $\hat{\mathbf{P}}_n$ and $\mathbf{Z}$. For any data generating process $P^o$ satisfying the
conditions of Exercise~4.26, we have that the asymptotic covariance matrix
of $\tilde{\bfbeta}_n$ is
\[
\mathbf{D}_n = (\mathbf{Q}_n'\mathbf{P}_n\mathbf{Q}_n)^{-1}\mathbf{Q}_n'\mathbf{P}_n\mathbf{V}_n\mathbf{P}_n\mathbf{Q}_n(\mathbf{Q}_n'\mathbf{P}_n\mathbf{Q}_n)^{-1}.
\]

We leave the dependence of $\mathbf{D}_n$ (and other quantities) on $P^o$ implicit for
now, and begin our analysis by considering how to choose $\hat{\mathbf{P}}_n$ to make $\mathbf{D}_n$
as small as possible.

Until now, we have let $\hat{\mathbf{P}}_n$ be any positive definite matrix. It turns out,
however, that by choosing $\hat{\mathbf{P}}_n = \hat{\mathbf{V}}_n^{-1}$, one obtains an asymptotically efficient estimator for the class of IV estimators with given instrumental
variables $\mathbf{Z}$. To prove this, we make use of the following proposition.

\begin{proposition}\label{prop:4.44}
Let $\mathbf{A}$ and $\mathbf{B}$ be positive definite matrices of order $k$. Then $\mathbf{A} - \mathbf{B}$ is positive semidefinite if and only if $\mathbf{B}^{-1} - \mathbf{A}^{-1}$ is positive semidefinite.
\end{proposition}

\begin{proof}
This follows from Goldberger (1964, Theorem~1.7.21, p.~38).
\end{proof}

This result is useful because in the cases of interest to us it will often be
much easier to determine whether $\mathbf{B}^{-1} - \mathbf{A}^{-1}$ is positive semidefinite than
to examine the positive semidefiniteness of $\mathbf{A} - \mathbf{B}$ directly.

\begin{proposition}\label{prop:4.45}
Given instrumental variables $\mathbf{Z}$, let $\mathcal{P}$ be the collection of probability measures satisfying the conditions of Exercise~4.26. Then the choice $\hat{\mathbf{P}}_n = \hat{\mathbf{V}}_n^{-1}$ gives the IV estimator
\[
\bfbeta_n^* = (\mathbf{X}'\mathbf{Z}\hat{\mathbf{V}}_n^{-1}\mathbf{Z}'\mathbf{X})^{-1}\mathbf{X}'\mathbf{Z}\hat{\mathbf{V}}_n^{-1}\mathbf{Z}'\mathbf{Y},
\]
which is asymptotically efficient for $\mathcal{P}$ within the class $\mathcal{E}$ of instrumental variables estimators of the form
\[
\tilde{\bfbeta}_n = (\mathbf{X}'\mathbf{Z}\hat{\mathbf{P}}_n\mathbf{Z}'\mathbf{X})^{-1}\mathbf{X}'\mathbf{Z}\hat{\mathbf{P}}_n\mathbf{Z}'\mathbf{Y}.
\]
\end{proposition}

\begin{proof}
From Exercise~4.26, we have
\begin{eqnarray*}
\mathrm{avar}(\tilde{\bfbeta}_n) &=& (\mathbf{Q}_n'\mathbf{P}_n\mathbf{Q}_n)^{-1}\mathbf{Q}_n'\mathbf{P}_n\mathbf{V}_n\mathbf{P}_n\mathbf{Q}_n(\mathbf{Q}_n'\mathbf{P}_n\mathbf{Q}_n)^{-1},\\
\mathrm{avar}(\bfbeta_n^*) &=& (\mathbf{Q}_n'\mathbf{V}_n^{-1}\mathbf{Q}_n)^{-1}.
\end{eqnarray*}
We apply Proposition~4.44 and consider
\begin{align*}
&(\mathrm{avar}(\bfbeta_n^*))^{-1} - (\mathrm{avar}(\tilde{\bfbeta}_n))^{-1} \\
&= \mathbf{Q}_n'\mathbf{V}_n^{-1}\mathbf{Q}_n - \mathbf{Q}_n'\mathbf{P}_n\mathbf{Q}_n(\mathbf{Q}_n'\mathbf{P}_n\mathbf{V}_n\mathbf{P}_n\mathbf{Q}_n)^{-1}\mathbf{Q}_n'\mathbf{P}_n\mathbf{Q}_n \\
&= \mathbf{Q}_n'\mathbf{V}_n^{-1/2}(\mathbf{I} \\
&\quad -\mathbf{V}_n^{1/2}\mathbf{P}_n\mathbf{Q}_n(\mathbf{Q}_n'\mathbf{P}_n\mathbf{V}_n^{1/2}\mathbf{V}_n^{1/2}\mathbf{P}_n\mathbf{Q}_n)^{-1}\mathbf{Q}_n'\mathbf{P}_n\mathbf{V}_n^{1/2})\mathbf{V}_n^{-1/2}\mathbf{Q}_n \\
&= \mathbf{Q}_n'\mathbf{V}_n^{-1/2}(\mathbf{I} - \mathbf{G}_n(\mathbf{G}_n'\mathbf{G}_n)^{-1}\mathbf{G}_n')\mathbf{V}_n^{-1/2}\mathbf{Q}_n,
\end{align*}
where $\mathbf{G}_n \equiv \mathbf{V}_n^{-1/2}\mathbf{P}_n\mathbf{Q}_n$. This is a quadratic form in an idempotent matrix
and is therefore p.s.d. As this holds for all $P^o$ in $\mathcal{P}$, the result follows.
\end{proof}

Hansen (1982) considers estimators for coefficients $\bfbeta_o$ implicitly defined
by moment conditions of the form $E(\mathbf{g}(\mathbf{X}_t, \mathbf{Y}_t, \mathbf{Z}_t, \bfbeta_o)) = \mathbf{0}$, where $\mathbf{g}$ is
an $l \times 1$ vector-valued function. Analogous to our construction of the IV
estimator, Hansen's estimator is constructed by attempting to make the
sample moment
\[
n^{-1}\sum_{t=1}^{n} \mathbf{g}(\mathbf{X}_t, \mathbf{Y}_t, \mathbf{Z}_t, \bfbeta)
\]
as close to zero as possible by solving the problem
\[
\min_{\bfbeta} \left[n^{-1}\sum_{t=1}^{n} \mathbf{g}(\mathbf{X}_t, \mathbf{Y}_t, \mathbf{Z}_t, \bfbeta)\right]'\hat{\mathbf{P}}_n\left[n^{-1}\sum_{t=1}^{n} \mathbf{g}(\mathbf{X}_t, \mathbf{Y}_t, \mathbf{Z}_t, \bfbeta)\right].
\]
Hansen establishes that the resulting ``method of moments'' estimator is
consistent and asymptotically normal and that the choice $\hat{\mathbf{P}}_n = \hat{\mathbf{V}}_n^{-1}$,
$\hat{\mathbf{V}}_n - \mathbf{V}_n \convp \mathbf{0}$, where
\[
\mathbf{V}_n = \mathrm{var}(n^{-1/2}\sum_{t=1}^{n} \mathbf{g}(\mathbf{X}_t, \mathbf{Y}_t, \mathbf{Z}_t, \bfbeta_o))
\]
delivers the asymptotically efficient estimator in the class of method of
moment estimators. Hansen calls the method of moments estimator with
$\hat{\mathbf{P}}_n = \hat{\mathbf{V}}_n^{-1}$ the \emph{generalized method of moments} (GMM) estimator.

Thus, the estimator $\bfbeta_n^*$ of Proposition~4.45 is the GMM estimator for
the case in which
\[
\mathbf{g}(\mathbf{X}_t, \mathbf{Y}_t, \mathbf{Z}_t, \bfbeta_o) = \mathbf{Z}_t(\mathbf{Y}_t - \mathbf{X}_t'\bfbeta_o).
\]
We call $\bfbeta_n^*$ a ``linear'' GMM estimator because the defining moment condition $E(\mathbf{Z}_t(\mathbf{Y}_t - \mathbf{X}_t'\bfbeta_o)) = \mathbf{0}$ is linear in $\bfbeta_o$.

\begin{exercises}
\item\label{ex:4.46} \textit{Given instrumental variables\/ $\mathbf{X}$, suppose that}
\[
\mathbf{V}_n^{-1/2}\sum \mathbf{X}_t\eps_t \overset{A}{\sim} N(\mathbf{0}, \mathbf{I}),
\]
\textit{where $\mathbf{V}_n = \sigma_o^2\mathbf{M}_n$. Show that the asymptotically efficient IV estimator is the least squares estimator, $\hat{\bfbeta}_n$, according to Proposition~4.45.}

\item\label{ex:4.47} \textit{Given instrumental variables\/ $\mathbf{Z}$, suppose that}
\[
\mathbf{V}_n^{-1/2}\sum \mathbf{Z}_t\eps_t \overset{A}{\sim} N(\mathbf{0}, \mathbf{I}),
\]
\textit{where $\mathbf{V}_n = \sigma_o^2\mathbf{L}_n$ and $\mathbf{L}_n \equiv E(\mathbf{Z}'\mathbf{Z}/n)$. Show that the asymptotically
efficient IV estimator is the two-stage least squares estimator}
\[
\tilde{\bfbeta}_{2SLS} = (\mathbf{X}'\mathbf{Z}(\mathbf{Z}'\mathbf{Z})^{-1}\mathbf{Z}'\mathbf{X})^{-1}\mathbf{X}'\mathbf{Z}(\mathbf{Z}'\mathbf{Z})^{-1}\mathbf{Z}'\mathbf{Y}
\]
\textit{according to Proposition~4.45.}

\item\label{ex:4.48} \textit{Define}
\begin{align*}
\bfbeta_n^{**} &= \bfbeta_n^* - (\mathbf{X}'\mathbf{Z}\hat{\mathbf{P}}_n\mathbf{Z}'\mathbf{X})^{-1}\nabla \mathbf{s}(\bfbeta_n^*) \\
&\quad\times[\nabla \mathbf{s}(\bfbeta_n^*)'(\mathbf{X}'\mathbf{Z}\hat{\mathbf{P}}_n\mathbf{Z}'\mathbf{X})^{-1}\nabla \mathbf{s}(\bfbeta_n^*)]^{-1}\mathbf{s}(\bfbeta_n^*).
\end{align*}
\textit{Show that under the conditions of Exercise~4.26 and the conditions on\/ $\mathbf{s}$
that $\sqrt{n}(\bfbeta_n^{**} - \bfbeta_n^*) \convp \mathbf{0}$, so that $\sqrt{n}(\bfbeta_n^* - \bfbeta_o)$ has the same asymptotic
distribution as $\sqrt{n}(\bfbeta_n^{**} - \bfbeta_o)$. (Hint: Show that $\sqrt{n}\mathbf{s}(\bfbeta_n^*) \convp \mathbf{0}$.)}
\end{exercises}

Note that the value of $\sigma_o^2$ plays no role in either Exercise~4.46 or Exercise
4.47 as long as it is finite. In what follows, we shall simply ignore $\sigma_o^2$, and
proceed as if $\sigma_o^2 = 1$.

If instead of testing the restrictions $\mathbf{s}(\bfbeta_o) = \mathbf{0}$, as considered in the previous section, we believe or know these restrictions to be true (because,
for example, they are dictated by economic theory), then, as we discuss
next, we can further improve asymptotic efficiency by imposing these restrictions. (Of course, for this to really work, the restrictions must in fact
be true.)

Thus, suppose we are given constraints $\mathbf{s}(\bfbeta_o) = \mathbf{0}$, where $\mathbf{s} : \mathbb{R}^k \to \mathbb{R}^q$ is
a known continuously differentiable function, such that $\mathrm{rank}(\nabla \mathbf{s}(\bfbeta_o)) = q$
and $\nabla \mathbf{s}(\bfbeta_o)$ ($k \times q$) is finite; The \emph{constrained instrumental variables estimator} can be found as the solution to the problem
\[
\min_{\bfbeta} (\mathbf{Y} - \mathbf{X}\bfbeta)'\mathbf{Z}\hat{\mathbf{P}}_n\mathbf{Z}'(\mathbf{Y} - \mathbf{X}\bfbeta), \quad \text{s.t.}\ \mathbf{s}(\bfbeta) = \mathbf{0},
\]
which is equivalent to finding the saddle point of the Lagrangian
\[
\mathcal{L} = (\mathbf{Y} - \mathbf{X}\bfbeta)'\mathbf{Z}\hat{\mathbf{P}}_n\mathbf{Z}'(\mathbf{Y} - \mathbf{X}\bfbeta) + \mathbf{s}(\bfbeta)'\boldsymbol{\lambda}.
\]

The first-order conditions are
\begin{eqnarray*}
\frac{\partial \mathcal{L}}{\partial \bfbeta} &=& 2(\mathbf{X}'\mathbf{Z}\hat{\mathbf{P}}_n\mathbf{Z}'\mathbf{X})\bfbeta - 2\mathbf{X}'\mathbf{Z}\hat{\mathbf{P}}_n\mathbf{Z}'\mathbf{Y} + \nabla \mathbf{s}(\bfbeta)\boldsymbol{\lambda} = \mathbf{0}\\
\frac{\partial \mathcal{L}}{\partial \boldsymbol{\lambda}} &=& \mathbf{s}(\bfbeta) = \mathbf{0}.
\end{eqnarray*}

Setting $\tilde{\bfbeta}_n = (\mathbf{X}'\mathbf{Z}\hat{\mathbf{P}}_n\mathbf{Z}'\mathbf{X})^{-1}\mathbf{X}'\mathbf{Z}\hat{\mathbf{P}}_n\mathbf{Z}'\mathbf{Y}$ and taking a mean value expansion of $\mathbf{s}(\tilde{\bfbeta})$ around $\mathbf{s}(\tilde{\bfbeta}_n)$ yields the equations
\begin{eqnarray*}
\frac{\partial \mathcal{L}}{\partial \bfbeta} &=& 2(\mathbf{X}'\mathbf{Z}\hat{\mathbf{P}}_n\mathbf{Z}'\mathbf{X})(\bfbeta - \tilde{\bfbeta}_n) + \nabla \mathbf{s}(\bfbeta)\boldsymbol{\lambda} = \mathbf{0},\\
\frac{\partial \mathcal{L}}{\partial \boldsymbol{\lambda}} &=& \mathbf{s}(\tilde{\bfbeta}_n) + \nabla \bar{\mathbf{s}}'(\bfbeta - \tilde{\bfbeta}_n) = \mathbf{0},
\end{eqnarray*}
where $\nabla \bar{\mathbf{s}}'$ is the $q \times k$ Jacobian matrix with $i$th row evaluated at a mean
value $\bar{\bfbeta}_n^{(i)}$. To solve for $\boldsymbol{\lambda}$ in the first equation, premultiply by $\nabla \bar{\mathbf{s}}'(\mathbf{X}'\mathbf{Z}\hat{\mathbf{P}}_n\mathbf{Z}'\mathbf{X})^{-1}$ to get
\[
2\nabla \bar{\mathbf{s}}(\bfbeta - \tilde{\bfbeta}_n) + \nabla \bar{\mathbf{s}}'(\mathbf{X}'\mathbf{Z}\hat{\mathbf{P}}_n\mathbf{Z}'\mathbf{X})^{-1}\nabla \mathbf{s}(\bfbeta)\boldsymbol{\lambda} = \mathbf{0},
\]
substitute $-\mathbf{s}(\tilde{\bfbeta}_n) = \nabla \bar{\mathbf{s}}'(\bfbeta - \tilde{\bfbeta}_n)$, and invert $\nabla \bar{\mathbf{s}}'(\mathbf{X}'\mathbf{Z}\hat{\mathbf{P}}_n\mathbf{Z}'\mathbf{X})^{-1}\nabla \mathbf{s}(\bfbeta)$ to
obtain
\[
\boldsymbol{\lambda} = 2[\nabla \bar{\mathbf{s}}'(\mathbf{X}'\mathbf{Z}\hat{\mathbf{P}}_n\mathbf{Z}'\mathbf{X})^{-1}\nabla \mathbf{s}(\bfbeta)]^{-1}\mathbf{s}(\tilde{\bfbeta}_n).
\]
The expression for $\partial \mathcal{L}/\partial \bfbeta$ above yields
\[
\bfbeta - \tilde{\bfbeta}_n = -(\mathbf{X}'\mathbf{Z}\hat{\mathbf{P}}_n\mathbf{Z}'\mathbf{X})^{-1}\nabla \mathbf{s}(\bfbeta)\boldsymbol{\lambda}/2
\]
so we obtain the solution for $\bfbeta$ by substituting for $\boldsymbol{\lambda}$:
\[
\bfbeta = \tilde{\bfbeta}_n - (\mathbf{X}'\mathbf{Z}\hat{\mathbf{P}}_n\mathbf{Z}'\mathbf{X})^{-1}\nabla \mathbf{s}(\bfbeta)[\nabla \bar{\mathbf{s}}'(\mathbf{X}'\mathbf{Z}\hat{\mathbf{P}}_n\mathbf{Z}'\mathbf{X})^{-1}\nabla \mathbf{s}(\bfbeta)]^{-1}\mathbf{s}(\tilde{\bfbeta}_n).
\]
The difficulty with this solution is that it is not in closed form, because the
unknown $\bfbeta$ appears on both sides of the equation. Further, appearing in
this expression is $\nabla \bar{\mathbf{s}}$, which has $q$ rows each of which depends on a mean
value lying between $\bfbeta$ and $\tilde{\bfbeta}_n$.

Nevertheless, a computationally practical and asymptotically equivalent
result can be obtained by replacing $\nabla \mathbf{s}(\bfbeta)$ and $\nabla \bar{\mathbf{s}}$ by $\nabla \mathbf{s}(\tilde{\bfbeta}_n)$ on the right-hand side of the expression above, which yields
\begin{align*}
\bfbeta_n^* &= \tilde{\bfbeta}_n - (\mathbf{X}'\mathbf{Z}\hat{\mathbf{P}}_n\mathbf{Z}'\mathbf{X})^{-1}\nabla \mathbf{s}(\tilde{\bfbeta}_n)[\nabla \mathbf{s}(\tilde{\bfbeta}_n)' \\
&\quad\times(\mathbf{X}'\mathbf{Z}\hat{\mathbf{P}}_n\mathbf{Z}'\mathbf{X})^{-1}\nabla \mathbf{s}(\tilde{\bfbeta}_n)]^{-1}\mathbf{s}(\tilde{\bfbeta}_n).
\end{align*}
This gives us a convenient way of computing a constrained IV estimator.
First we compute the unconstrained estimator, and then we ``impose'' the
constraints by subtracting a ``correction factor''
\[
(\mathbf{X}'\mathbf{Z}\hat{\mathbf{P}}_n\mathbf{Z}'\mathbf{X})^{-1}\nabla \mathbf{s}(\tilde{\bfbeta}_n)[\nabla \mathbf{s}(\tilde{\bfbeta}_n)'(\mathbf{X}'\mathbf{Z}\hat{\mathbf{P}}_n\mathbf{Z}'\mathbf{X})^{-1}\nabla \mathbf{s}(\tilde{\bfbeta}_n)]^{-1}\mathbf{s}(\tilde{\bfbeta}_n)
\]
from the unconstrained estimator. We say ``impose'' because $\bfbeta_n^*$ will not
satisfy the constraints exactly for any finite $n$. However, an estimator that
does satisfy the constraints to any desired degree of accuracy can be obtained by iterating the procedure just described, that is, by replacing $\tilde{\bfbeta}_n$
by $\bfbeta_n^*$ in the formula above to get a second round estimator, say, $\bfbeta_n^{**}$.
This process could continue until the change in the resulting estimator was
sufficiently small. Nevertheless, this iteration process has no effect on the
asymptotic covariance matrix of the resulting estimator.

Thus, in considering the effect of imposing the restriction $\mathbf{s}(\bfbeta_o) = \mathbf{0}$, we
can just compare the asymptotic behavior of $\bfbeta_n^*$ to $\tilde{\bfbeta}_n$.

As we saw in Proposition~4.45, the asymptotically efficient IV estimator
takes $\hat{\mathbf{P}}_n = \hat{\mathbf{V}}_n^{-1}$. We therefore consider the effect of imposing constraints
on the IV estimator with $\hat{\mathbf{P}}_n = \hat{\mathbf{V}}_n^{-1}$.

\begin{theorem}\label{thm:4.49}
Suppose the conditions of Exercise~4.26 hold for $\hat{\mathbf{P}}_n = \hat{\mathbf{V}}_n^{-1}$ and that $\mathbf{s} : \mathbb{R}^k \to \mathbb{R}^q$ is a continuously differentiable function such that $\mathbf{s}(\bfbeta_o) = 0$, $\nabla \mathbf{s}(\bfbeta_o)$ is finite and $\mathrm{rank}(\nabla \mathbf{s}(\bfbeta_o)) = q$. Define $\tilde{\bfbeta}_n = (\mathbf{X}'\mathbf{Z}\hat{\mathbf{V}}_n^{-1}\mathbf{Z}'\mathbf{X})^{-1}\mathbf{X}'\mathbf{Z}\hat{\mathbf{V}}_n^{-1}\mathbf{Z}'\mathbf{Y}$ and define $\bfbeta_n^*$ as the constrained IV estimator above with $\hat{\mathbf{P}}_n = \hat{\mathbf{V}}_n^{-1}$. Then
\begin{align*}
\mathrm{avar}(\tilde{\bfbeta}_n) - \mathrm{avar}(\bfbeta_n^*) &= \mathrm{avar}(\tilde{\bfbeta}_n)\nabla \mathbf{s}(\bfbeta_o)[\nabla \mathbf{s}(\bfbeta_o)'\mathrm{avar}(\tilde{\bfbeta}_n)\nabla \mathbf{s}(\bfbeta_o)]^{-1} \\
&\quad\times \nabla \mathbf{s}(\bfbeta_o)'\mathrm{avar}(\tilde{\bfbeta}_n),
\end{align*}
which is a positive semidefinite matrix.
\end{theorem}

\begin{proof}
From Exercise~4.26 it follows that
\[
\mathrm{avar}(\tilde{\bfbeta}_n) = (\mathbf{Q}_n'\mathbf{V}_n^{-1}\mathbf{Q}_n)^{-1}.
\]
Taking a mean value expansion of $\mathbf{s}(\tilde{\bfbeta}_n)$ around $\bfbeta_o$ gives $\mathbf{s}(\tilde{\bfbeta}_n) = \mathbf{s}(\bfbeta_o) +
\nabla \bar{\mathbf{s}}'(\tilde{\bfbeta}_n - \bfbeta_o)$, and because $\mathbf{s}(\bfbeta_o) = \mathbf{0}$, we have $\mathbf{s}(\tilde{\bfbeta}_n) = \nabla \bar{\mathbf{s}}'(\tilde{\bfbeta}_n - \bfbeta_o)$.
Substituting this in the formula for $\bfbeta_n^*$ allows us to write
\[
\sqrt{n}(\bfbeta_n^* - \bfbeta_o) = \tilde{\mathbf{A}}_n\sqrt{n}(\tilde{\bfbeta}_n - \bfbeta_o),
\]
where
\begin{align*}
\tilde{\mathbf{A}}_n &= \mathbf{I} - (\mathbf{X}'\mathbf{Z}\hat{\mathbf{V}}_n^{-1}\mathbf{Z}'\mathbf{X})^{-1}\nabla \mathbf{s}(\tilde{\bfbeta}_n)[\nabla \mathbf{s}(\tilde{\bfbeta}_n)' \\
&\quad\times(\mathbf{X}'\mathbf{Z}\hat{\mathbf{V}}_n^{-1}\mathbf{Z}'\mathbf{X})^{-1}\nabla \mathbf{s}(\tilde{\bfbeta}_n)]^{-1}\nabla \bar{\mathbf{s}}'.
\end{align*}
Under the conditions of Exercise~4.26, $\tilde{\bfbeta}_n \convp \bfbeta_o$ and Proposition~2.30
applies to ensure that $\tilde{\mathbf{A}}_n - \mathbf{A}_n \convp \mathbf{0}$, where
\[
\mathbf{A}_n = \mathbf{I} - \mathrm{avar}(\tilde{\bfbeta}_n)\nabla \mathbf{s}(\bfbeta_o)[\nabla \mathbf{s}(\bfbeta_o)'\mathrm{avar}(\tilde{\bfbeta}_n)\nabla \mathbf{s}(\bfbeta_o)]^{-1}\nabla \mathbf{s}(\bfbeta_o)'.
\]
Hence, by Lemma~4.6,
\[
\sqrt{n}(\bfbeta_n^* - \bfbeta_o) - \mathbf{A}_n\sqrt{n}(\tilde{\bfbeta}_n - \bfbeta_o) = (\tilde{\mathbf{A}}_n - \mathbf{A}_n)\sqrt{n}(\tilde{\bfbeta}_n - \bfbeta_o) \convp \mathbf{0},
\]
because $\sqrt{n}(\tilde{\bfbeta}_n - \bfbeta_o)$ is $O_p(1)$ as a consequence of Exercise~4.26. From the
asymptotic equivalence Lemma~4.7 it follows that $\sqrt{n}(\bfbeta_n^* - \bfbeta_o)$ has the
same asymptotic distribution as $\mathbf{A}_n\sqrt{n}(\tilde{\bfbeta}_n - \bfbeta_o)$. It follows from Lemma
4.23 that $\sqrt{n}(\bfbeta_n^* - \bfbeta_o)$ is asymptotically normal with mean zero and
\[
\mathrm{avar}(\bfbeta_n^*) = \boldsymbol{\Gamma}_n = \mathbf{A}_n\mathrm{avar}(\tilde{\bfbeta}_n)\mathbf{A}_n'.
\]
Straightforward algebra yields
\begin{align*}
\mathrm{avar}(\bfbeta_n^*) &= \mathrm{avar}(\tilde{\bfbeta}_n) - \mathrm{avar}(\tilde{\bfbeta}_n)\nabla \mathbf{s}(\bfbeta_o) \\
&\quad\times[\nabla \mathbf{s}(\bfbeta_o)'\mathrm{avar}(\tilde{\bfbeta}_n)\nabla \mathbf{s}(\bfbeta_o)]^{-1}\nabla \mathbf{s}(\bfbeta_o)'\mathrm{avar}(\tilde{\bfbeta}_n),
\end{align*}
and the result follows immediately.
\end{proof}

This result guarantees that imposing correct a priori restrictions leads
to an efficiency improvement over the efficient IV estimator that does not
impose these restrictions. Interestingly, imposing the restrictions using the
formula for $\bfbeta_n^*$ with an inefficient estimator $\tilde{\bfbeta}_n$ for given instrumental variables may or may not lead to efficiency gains relative to $\tilde{\bfbeta}_n$.

A feature of this result worth noting is that the asymptotic distribution
of the constrained estimator $\bfbeta_n^*$ is concentrated in a $k - q$-dimensional
subspace of $\mathbb{R}^k$, so that $\mathrm{avar}(\bfbeta_n^*)$ will not be nonsingular. Instead, it has
rank $k - q$. In particular, observe that
\[
\sqrt{n}(\bfbeta_n^* - \bfbeta_o) = \sqrt{n}\mathbf{A}_n(\tilde{\bfbeta}_n - \bfbeta_o).
\]
But
\begin{align*}
\nabla \mathbf{s}(\bfbeta_o)'\mathbf{A}_n &= \nabla \mathbf{s}(\bfbeta_o)'\left(\mathbf{I} - \mathrm{avar}(\tilde{\bfbeta}_n)\nabla \mathbf{s}(\bfbeta_o)\right. \\
&\quad\left.\times[\nabla \mathbf{s}(\bfbeta_o)'\mathrm{avar}(\tilde{\bfbeta}_n)\nabla \mathbf{s}(\bfbeta_o)]^{-1}\nabla \mathbf{s}(\bfbeta_o)'\right) \\
&= \nabla \mathbf{s}(\bfbeta_o)' - \nabla \mathbf{s}(\bfbeta_o)' \\
&= \mathbf{0}.
\end{align*}
Consequently, we have
\[
\mathrm{avar}(\sqrt{n}\,\nabla \mathbf{s}(\bfbeta_o)'(\bfbeta_n^* - \bfbeta_o)) = \mathbf{0}.
\]

An alternative way to improve the asymptotic efficiency of the IV estimator is to include additional valid instrumental variables. To establish
this, we make use of the formula for the inverse of a partitioned matrix,
which we now state for convenience.

\begin{proposition}\label{prop:4.50}
Define the $k \times k$ nonsingular symmetric matrix
\[
\mathbf{A} = \begin{bmatrix} \mathbf{B} & \mathbf{C}' \\ \mathbf{C} & \mathbf{D} \end{bmatrix},
\]
where $\mathbf{B}$ is $k_1 \times k_1$, $\mathbf{C}$ is $k_2 \times k_1$ and $\mathbf{D}$ is $k_2 \times k_2$. Then, defining $\mathbf{E} \equiv \mathbf{D} - \mathbf{C}\mathbf{B}^{-1}\mathbf{C}'$,
\[
\mathbf{A}^{-1} = \begin{bmatrix} \mathbf{B}^{-1}(\mathbf{I} + \mathbf{C}'\mathbf{E}^{-1}\mathbf{C}\mathbf{B}^{-1}) & -\mathbf{B}^{-1}\mathbf{C}'\mathbf{E}^{-1} \\ -\mathbf{E}^{-1}\mathbf{C}\mathbf{B}^{-1} & \mathbf{E}^{-1} \end{bmatrix}.
\]
\end{proposition}

\begin{proof}
See Goldberger (1964, p.~27).
\end{proof}

\begin{proposition}\label{prop:4.51}
Partition $\mathbf{Z}$ as $\mathbf{Z} = (\mathbf{Z}_1, \mathbf{Z}_2)$ and let $\mathcal{P}$ be the collection of all probability measures such that the conditions of Exercise~4.26 hold for both $\mathbf{Z}_1$ and $\mathbf{Z}_2$. Define $\mathbf{V}_{1n} = E(\mathbf{Z}_1'\bfeps\bfeps'\mathbf{Z}_1/n)$, and the estimators
\begin{eqnarray*}
\tilde{\bfbeta}_n &=& (\mathbf{X}'\mathbf{Z}_1\hat{\mathbf{V}}_{1n}^{-1}\mathbf{Z}_1'\mathbf{X})^{-1}\mathbf{X}'\mathbf{Z}_1\hat{\mathbf{V}}_{1n}^{-1}\mathbf{Z}_1'\mathbf{Y},\\
\bfbeta_n^* &=& (\mathbf{X}'\mathbf{Z}\hat{\mathbf{V}}_n^{-1}\mathbf{Z}'\mathbf{X})^{-1}\mathbf{X}'\mathbf{Z}\hat{\mathbf{V}}_n^{-1}\mathbf{Z}'\mathbf{Y}.
\end{eqnarray*}
Then for each $P^o$ in $\mathcal{P}$, $\mathrm{avar}(\tilde{\bfbeta}_n) - \mathrm{avar}(\bfbeta_n^*)$ is a positive semidefinite matrix for all $n$ sufficiently large.
\end{proposition}

\begin{proof}
Partition $\mathbf{Q}_n$ as $\mathbf{Q}_n' = (\mathbf{Q}_{1n}', \mathbf{Q}_{2n}')$, where $\mathbf{Q}_{1n} = E(\mathbf{Z}_1'\mathbf{X}/n)$, $\mathbf{Q}_{2n} = E(\mathbf{Z}_2'\mathbf{X}/n)$, and partition $\mathbf{V}_n$ as
\[
\mathbf{V}_n = \begin{bmatrix} \mathbf{V}_{1n} & \mathbf{V}_{12n} \\ \mathbf{V}_{21n} & \mathbf{V}_{2n} \end{bmatrix}.
\]
The partitioned inverse formula gives
\[
\mathbf{V}_n^{-1} = \begin{bmatrix} \mathbf{V}_{1n}^{-1}(\mathbf{I} + \mathbf{V}_{12n}\mathbf{E}_n^{-1}\mathbf{V}_{21n}\mathbf{V}_{1n}^{-1}) & -\mathbf{V}_{1n}^{-1}\mathbf{V}_{12n}\mathbf{E}_n^{-1} \\ -\mathbf{E}_n^{-1}\mathbf{V}_{21n}\mathbf{V}_{1n}^{-1} & \mathbf{E}_n^{-1} \end{bmatrix},
\]
where $\mathbf{E}_n = \mathbf{V}_{2n} - \mathbf{V}_{21n}\mathbf{V}_{1n}^{-1}\mathbf{V}_{12n}$. From Exercise~4.26, we have
\begin{eqnarray*}
\mathrm{avar}(\tilde{\bfbeta}_n) &=& (\mathbf{Q}_{1n}'\mathbf{V}_{1n}^{-1}\mathbf{Q}_{1n})^{-1},\\
\mathrm{avar}(\bfbeta_n^*) &=& (\mathbf{Q}_n'\mathbf{V}_n^{-1}\mathbf{Q}_n)^{-1}.
\end{eqnarray*}
We apply Proposition~4.44 and consider
\begin{align*}
&(\mathrm{avar}(\bfbeta_n^*))^{-1} - (\mathrm{avar}(\tilde{\bfbeta}_n))^{-1} \\
&= \mathbf{Q}_n'\mathbf{V}_n^{-1}\mathbf{Q}_n - \mathbf{Q}_{1n}'\mathbf{V}_{1n}^{-1}\mathbf{Q}_{1n} \\
&= \mathbf{Q}_{1n}'\left(\mathbf{V}_{1n}^{-1} + \mathbf{V}_{1n}^{-1}\mathbf{V}_{12n}\mathbf{E}_n^{-1}\mathbf{V}_{21n}\mathbf{V}_{1n}^{-1}\right)\mathbf{Q}_{1n} \\
&\quad - \mathbf{Q}_{2n}'\mathbf{E}_n^{-1}\mathbf{V}_{21n}\mathbf{V}_{1n}^{-1}\mathbf{Q}_{1n} - \mathbf{Q}_{1n}'\mathbf{V}_{1n}^{-1}\mathbf{V}_{12n}\mathbf{E}_n^{-1}\mathbf{Q}_{2n} \\
&\quad + \mathbf{Q}_{2n}'\mathbf{E}_n^{-1}\mathbf{Q}_{2n} - \mathbf{Q}_{1n}'\mathbf{V}_{1n}^{-1}\mathbf{Q}_{1n} \\
&= \left(\mathbf{Q}_{1n}'\mathbf{V}_{1n}^{-1}\mathbf{V}_{12n} - \mathbf{Q}_{2n}'\right)\mathbf{E}_n^{-1}\left(\mathbf{V}_{21n}\mathbf{V}_{1n}^{-1}\mathbf{Q}_{1n} - \mathbf{Q}_{2n}\right).
\end{align*}
Because $\mathbf{E}_n^{-1}$ is a symmetric positive definite matrix (why?), we can write
$\mathbf{E}_n^{-1} = \mathbf{E}_n^{-1/2}\mathbf{E}_n^{-1/2}$, so that
\begin{align*}
&\left(\mathrm{avar}(\bfbeta_n^*)\right)^{-1} - \left(\mathrm{avar}(\tilde{\bfbeta}_n)\right)^{-1} \\
&= \left(\mathbf{Q}_{1n}'\mathbf{V}_{1n}^{-1}\mathbf{V}_{12n} - \mathbf{Q}_{2n}'\right)\mathbf{E}_n^{-1/2}\mathbf{E}_n^{-1/2}\left(\mathbf{V}_{21n}\mathbf{V}_{1n}^{-1}\mathbf{Q}_{1n} - \mathbf{Q}_{2n}\right).
\end{align*}
Because this is the product of a matrix and its transpose, we immediately
have $(\mathrm{avar}(\bfbeta_n^*))^{-1} - (\mathrm{avar}(\tilde{\bfbeta}_n))^{-1}$ is p.s.d. As this holds for all $P^o$ in $\mathcal{P}$ the
result follows, which by Proposition~4.44 implies that $\mathrm{avar}(\tilde{\bfbeta}_n) - \mathrm{avar}(\bfbeta_n^*)$
is p.s.d.
\end{proof}

This result states essentially that the asymptotic precision of the IV
estimator cannot be worsened by including additional instrumental variables. We can be more specific, however, and specify situations in which
$\mathrm{avar}(\bfbeta_n^*) = \mathrm{avar}(\tilde{\bfbeta}_n)$, so that nothing is gained by adding an extra instrumental variable.

\begin{proposition}\label{prop:4.52}
Let the conditions of Proposition~4.51 hold. Then we have $\mathrm{avar}(\tilde{\bfbeta}_n) = \mathrm{avar}(\bfbeta_n^*)$ if and only if
\[
E(\mathbf{X}'\mathbf{Z}_1/n)E(\mathbf{Z}_1'\bfeps\bfeps'\mathbf{Z}_1/n)^{-1}E(\mathbf{Z}_1'\bfeps\bfeps'\mathbf{Z}_2/n) - E(\mathbf{X}'\mathbf{Z}_2/n) = \mathbf{0}.
\]
\end{proposition}

\begin{proof}
Immediate from the final line of the proof of Proposition~4.51.
\end{proof}

To interpret this condition, consider the special case in which
\[
E(\mathbf{Z}'\bfeps\bfeps'\mathbf{Z}/n) = E(\mathbf{Z}'\mathbf{Z}/n).
\]
In this case the difference in Proposition~4.52 can be consistently estimated
by
\[
n^{-1}(\mathbf{X}'\mathbf{Z}_1(\mathbf{Z}_1'\mathbf{Z}_1)^{-1}\mathbf{Z}_1'\mathbf{Z}_2 - \mathbf{X}'\mathbf{Z}_2) = n^{-1}\mathbf{X}\left[\mathbf{Z}_1(\mathbf{Z}_1'\mathbf{Z}_1)^{-1}\mathbf{Z}_1' - \mathbf{I}\right]\mathbf{Z}_2.
\]
This quantity is recognizable as the cross product of $\mathbf{Z}_2$ and the projection
of $\mathbf{X}$ onto the space orthogonal to that spanned by the columns of $\mathbf{Z}_1$. If
we write $\tilde{\mathbf{X}} = \mathbf{X}(\mathbf{Z}_1'\mathbf{Z}_1)^{-1}\mathbf{Z}_1' - \mathbf{I})$, the difference in Proposition~4.52 is
consistently estimated by $\tilde{\mathbf{X}}'\mathbf{Z}_2/n$, so that $\mathrm{avar}(\tilde{\bfbeta}_n) = \mathrm{avar}(\bfbeta_n^*)$ if and only
if $\tilde{\mathbf{X}}'\mathbf{Z}_2/n \convp \mathbf{0}$, which can be interpreted as saying that adding $\mathbf{Z}_2$ to the
list of instrumental variables is of no use if it is uncorrelated with $\tilde{\mathbf{X}}$, the
matrix of residuals of the regression of $\mathbf{X}$ on $\mathbf{Z}_1$.

One of the interesting consequences of Propositions~4.51 and 4.52 is that
in the presence of heteroskedasticity or serial correlation of unknown form,
there may exist estimators for the linear model more efficient than OLS.
This result has been obtained independently by Cragg (1983) and Chamberlain (1982). To construct these estimators, it is necessary to find additional instrumental variables uncorrelated with $\eps_t$. If $E(\eps_t|\mathbf{X}_t) = \mathbf{0}$, such
instrumental variables are easily found because any measurable function
of $\mathbf{X}_t$ will be uncorrelated with $\eps_t$. Hence, we can set $\mathbf{Z}_t = (\mathbf{X}_t', \mathbf{z}(\mathbf{X}_t)')'$,
where $\mathbf{z}(\mathbf{X}_t)$ is a $(l - k) \times 1$ vector of measurable functions of $\mathbf{X}_t$.

\begin{example}\label{ex:4.53}
Let $p = k = 1$ so $Y_t = X_t\beta_o + \eps_t$, where $Y_t$, $X_t$ and $\eps_t$
are scalars. Suppose that $X_t$ is nonstochastic and for convenience suppose
$M_n = n^{-1}\sum X_t^2 \to 1$. Let $\eps_t$ be independent heterogeneously distributed
such that $E(\eps_t) = 0$ and $E(\eps_t^2) = \sigma_t^2$. Further, suppose $X_t > \delta > 0$
for all $t$, and take $z(X_t) = X_t^{-1}$ so that $\mathbf{Z}_t = (X_t, X_t^{-1})'$. We consider
$\hat{\beta}_n = (\mathbf{X}'\mathbf{X})^{-1}\mathbf{X}'\mathbf{Y}$ and $\beta_n^* = (\mathbf{X}'\mathbf{Z}\hat{\mathbf{V}}_n^{-1}\mathbf{Z}'\mathbf{X})^{-1}\mathbf{X}'\mathbf{Z}\hat{\mathbf{V}}_n^{-1}\mathbf{Z}'\mathbf{Y}$, and suppose
that sufficient other assumptions guarantee that the result of Exercise~4.26
holds for both estimators. By Propositions~4.51 and 4.52, it follows that
$\mathrm{avar}(\hat{\beta}_n) > \mathrm{avar}(\beta_n^*)$ if and only if
\[
\left(n^{-1}\sum_{t=1}^{n}\sigma_t^2 X_t^2\right)^{-1}\left(n^{-1}\sum_{t=1}^{n}\sigma_t^2\right) - 1 \neq 0
\]
or equivalently, if and only if $n^{-1}\sum_{t=1}^{n}\sigma_t^2 X_t^2 \neq n^{-1}\sum_{t=1}^{n}\sigma_t^2$. This would
certainly occur if $\sigma_t = X_t^{-1}$. (Verify this using Jensen's inequality.)
\end{example}

It also follows from Propositions~4.51 and 4.52 that when $\mathbf{V}_n \neq \mathbf{L}_n$, there
may exist estimators more efficient than two-stage least squares. If $l > k$,
additional instrumental variables are not necessarily required to improve
efficiency over 2SLS (see White, 1982); but as the result of Proposition
4.51 indicates, additional instrumental variables (e.g., functions of $\mathbf{Z}_t$) can
nevertheless generate further improvements.

The results so far suggest that in particular situations efficiency may or
may not be improved by using additional instrumental variables. This leads
us to ask whether there exists an optimal set of instrumental variables, that
is, a set of instrumental variables that one could not improve upon by using
additional instrumental variables. If so, we could in principle definitively
know whether or not a given set of instrumental variables is best.

To address this question, we introduce an approach developed by Bates
and White (1993) that permits direct construction of the optimal instrumental variables. This is in sharp contrast to the approach taken when
we considered the best choice for $\hat{\mathbf{P}}_n$. There, we seemingly pulled the best
choice, $\hat{\mathbf{V}}_n^{-1}$, out of a hat (as if by magic!), and simply verified its optimality. There was no hint of how we arrived at this choice. (This is a common
feature of the efficiency literature, a feature that has endowed it with no
little mystery.) With the Bates--White method, however, we will be able
to see how to construct the optimal (most efficient) instrumental variables
estimators, step-by-step.

Bates and White consider classes of estimators $\mathcal{E}$ indexed by a parameter
$\boldsymbol{\gamma}$ such that for each data generating process $P^o$ in a set $\mathcal{P}$,
\begin{equation}\label{eq:4.1}
\sqrt{n}(\hat{\boldsymbol{\theta}}_n(\boldsymbol{\gamma}) - \boldsymbol{\theta}^o) = \mathbf{H}_n^o(\boldsymbol{\gamma})^{-1}\sqrt{n}\mathbf{s}_n^o(\boldsymbol{\gamma}) + o_{p^o}(1),
\end{equation}
where: $\hat{\boldsymbol{\theta}}_n(\boldsymbol{\gamma})$ is an estimator indexed by $\boldsymbol{\gamma}$ taking values in a set $\Gamma$; $\boldsymbol{\theta}^o = \boldsymbol{\theta}(P^o)$ is the $k \times 1$ parameter vector corresponding to the data generating
process $P^o$; $\mathbf{H}_n^o(\boldsymbol{\gamma})$ is a nonstochastic nonsingular $k \times k$ matrix depending
on $\boldsymbol{\gamma}$ and $P^o$; $\mathbf{s}_n^o(\boldsymbol{\gamma})$ is a random $k \times 1$ vector depending on $\boldsymbol{\gamma}$ and $P^o$ such
that $\mathbf{I}_n^o(\boldsymbol{\gamma})^{-1/2}\sqrt{n}\mathbf{s}_n^o(\boldsymbol{\gamma}) \xrightarrow{d^o} N(\mathbf{0}, \mathbf{I})$, where $\mathbf{I}_n^o(\boldsymbol{\gamma}) \equiv \mathrm{var}^o(\sqrt{n}\mathbf{s}_n^o(\boldsymbol{\gamma}))$, $\xrightarrow{d^o}$
denotes convergence in distribution under $P^o$, and $\mathrm{var}^o$ denotes variance
under $P^o$; and $o_{p^o}(1)$ denotes terms vanishing in probability-$P^o$. Bates and
White consider how to choose $\boldsymbol{\gamma}$ in $\Gamma$ to obtain an asymptotically efficient
estimator.

Bates and White's method is actually formulated for a more general
setting but \eqref{eq:4.1} will suffice for us here.

To relate the Bates--White method to our instrumental variables estimation framework, we begin by writing the process generating the dependent
variables in a way that makes explicit the role of $P^o$. Specifically, we write
\[
\mathbf{Y}_t^o = \mathbf{X}_t^{o\prime}\bfbeta^o + \eps_t, \quad t = 1, 2, \ldots\;.
\]
The data generating process $P^o$ governs
the probabilistic behavior of all of the random variables (measurable mappings from the underlying measurable space $(\Omega, \mathcal{F})$) in our system. Given
the way our data are generated, the probabilistic behavior of the mapping
$\eps_t$ is governed by $P^o$, but the mapping itself is not determined by $P^o$. We
thus do not give a ${}^o$ superscript to $\eps_t$.

On the other hand, different values $\bfbeta^o$ (corresponding to $\boldsymbol{\theta}^o$ in the Bates--White setting) result in different mappings for the dependent variables; we
acknowledge this by writing $\mathbf{Y}_t^o$, and we view $\bfbeta^o$ as determined by $P^o$,
i.e., $\bfbeta^o = \bfbeta(P^o)$. Further, because lags of $\mathbf{Y}_t^o$ or elements of $\mathbf{Y}_t^o$ itself
(recall the simultaneous equations framework) may appear as elements of
the explanatory variable matrix, different values for $\bfbeta^o$ may also result in
different mappings for the explanatory variables. We acknowledge this by
writing $\mathbf{X}_t^o$.

Next, consider the instrumental variables. These also are measurable
mappings from $\{\Omega, \mathcal{F}\}$ whose probabilistic behavior is governed by $P^o$. As
we will see, it is useful to permit these mappings to depend on $P^o$ as well.
Consequently, we write $\mathbf{Z}_t^o$. Because our goal is the choice of optimal instrumental variables, we let $\mathbf{Z}_t^o$ depend on $\boldsymbol{\gamma}$ and write $\mathbf{Z}_t^o(\boldsymbol{\gamma})$ to emphasize this
dependence. The optimal choice of $\boldsymbol{\gamma}$ will then deliver the optimal choice of
instrumental variables. Later, we will revisit and extend this interpretation
of $\boldsymbol{\gamma}$.

In Exercise~4.26 $(ii)$, we assume $\mathbf{V}_n^{-1/2}n^{-1/2}\mathbf{Z}'\bfeps \convd N(\mathbf{0}, \mathbf{I})$, where
$\mathbf{V}_n \equiv \mathrm{var}(n^{-1/2}\mathbf{Z}'\bfeps)$. For the Bates--White framework with $\mathbf{Z}_t^o(\boldsymbol{\gamma})$, this
becomes $\mathbf{V}_n^o(\boldsymbol{\gamma})^{-1/2}n^{-1/2}\mathbf{Z}^o(\boldsymbol{\gamma})'\bfeps \xrightarrow{d^o} N(\mathbf{0}, \mathbf{I})$, where $\mathbf{V}_n^o(\boldsymbol{\gamma}) \equiv \mathrm{var}^o(n^{-1/2}
\mathbf{Z}^o(\boldsymbol{\gamma})'\bfeps)$, and $\mathbf{Z}^o(\boldsymbol{\gamma})$ is the $np \times l$ matrix with blocks $\mathbf{Z}_t^o(\boldsymbol{\gamma})'$. In $(iii.a)$ we
assume $\mathbf{Z}'\mathbf{X}/n - \mathbf{Q}_n = o_p(1)$, where $\mathbf{Q}_n \equiv E(\mathbf{Z}'\mathbf{X}/n)$ is $O(1)$ with uniformly full column rank. This becomes $\mathbf{Z}^o(\boldsymbol{\gamma})'\mathbf{X}^o/n - \mathbf{Q}_n^o = o_{p^o}(1)$, where
$\mathbf{Q}_n^o \equiv E^o(\mathbf{Z}^o(\boldsymbol{\gamma})'\mathbf{X}^o/n)$ is $O(1)$ with uniformly full column rank; $E^o(\cdot)$ denotes expectation with respect to $P^o$. In $(iii.b)$ we assume $\hat{\mathbf{P}}_n - \mathbf{P}_n = o_p(1)$, where $\mathbf{P}_n$ is $O(1)$ and uniformly positive definite. This becomes
$\hat{\mathbf{P}}_n(\boldsymbol{\gamma}) - \mathbf{P}_n^o(\boldsymbol{\gamma}) = o_{p^o}(1)$, where $\mathbf{P}_n^o(\boldsymbol{\gamma})$ is $O(1)$ and uniformly positive definite. Finally, the instrumental variables estimator has the form
\[
\tilde{\bfbeta}_n(\boldsymbol{\gamma}) = \left(\mathbf{X}^{o\prime}\mathbf{Z}^o(\boldsymbol{\gamma})\hat{\mathbf{P}}_n(\boldsymbol{\gamma})\mathbf{Z}^o(\boldsymbol{\gamma})'\mathbf{X}^o\right)^{-1}\mathbf{X}^{o\prime}\mathbf{Z}^o(\boldsymbol{\gamma})\hat{\mathbf{P}}_n(\boldsymbol{\gamma})\mathbf{Z}^o(\boldsymbol{\gamma})'\mathbf{Y}^o.
\]
We will use the following further extension of Exercise~4.26.

\begin{theorem}\label{thm:4.54}
Let $\mathcal{P}$ be a set of probability measures $P^o$ on $(\Omega, \mathcal{F})$, and let $\Gamma$ be a set with elements $\boldsymbol{\gamma}$. Suppose that for each $P^o$ in $\mathcal{P}$
\begin{enumerate}[label=(\roman*)]
\item $\mathbf{Y}_t^o = \mathbf{X}_t^{o\prime}\bfbeta^o + \eps_t$, \quad $t = 1, 2, \ldots$\,, $\bfbeta^o \in \mathbb{R}^k$.

\emph{For each $\boldsymbol{\gamma}$ in $\Gamma$ and each $P^o$ in $\mathcal{P}$, let $\{\mathbf{Z}_{nt}^o(\boldsymbol{\gamma})\}$ be a double array of random $p \times l$ matrices, such that}

\item $n^{-1/2}\sum_{t=1}^{n}\mathbf{Z}_{nt}^o(\boldsymbol{\gamma})\eps_t - n^{1/2}\mathbf{m}_n^o(\boldsymbol{\gamma}) = o_{p^o}(1)$, where
\[
\mathbf{V}_n^o(\boldsymbol{\gamma})^{-1/2}n^{1/2}\mathbf{m}_n^o(\boldsymbol{\gamma}) \xrightarrow{d^o} N(\mathbf{0}, \mathbf{I}),
\]
and $\mathbf{V}_n^o(\boldsymbol{\gamma}) \equiv \mathrm{var}^o(n^{1/2}\mathbf{m}_n^o(\boldsymbol{\gamma}))$ is $O(1)$ and uniformly positive definite;

\item \begin{enumerate}[label=(\alph*)]
\item $n^{-1}\sum_{t=1}^{n}\mathbf{Z}_{nt}^o(\boldsymbol{\gamma})\mathbf{X}_t^{o\prime} - \mathbf{Q}_n^o(\boldsymbol{\gamma}) = o_{p^o}(1)$, where $\mathbf{Q}_n^o(\boldsymbol{\gamma})$ is $O(1)$ with uniformly full column rank;
\item There exists $\hat{\mathbf{P}}_n(\boldsymbol{\gamma})$ such that $\hat{\mathbf{P}}_n(\boldsymbol{\gamma}) - \mathbf{P}_n^o(\boldsymbol{\gamma}) = o_{p^o}(1)$ and $\mathbf{P}_n^o(\boldsymbol{\gamma}) = O(1)$ and is symmetric and uniformly positive definite.
\end{enumerate}
\end{enumerate}

\emph{Then for each $P^o$ in $\mathcal{P}$ and $\boldsymbol{\gamma}$ in $\Gamma$}
\[
\mathbf{D}_n^o(\boldsymbol{\gamma})^{-1/2}\sqrt{n}\left(\tilde{\bfbeta}_n(\boldsymbol{\gamma}) - \bfbeta^o\right) \xrightarrow{d^o} N(\mathbf{0}, \mathbf{I}),
\]
\emph{where}
\begin{align*}
\mathbf{D}_n^o(\boldsymbol{\gamma}) &\equiv (\mathbf{Q}_n^o(\boldsymbol{\gamma})'\mathbf{P}_n^o(\boldsymbol{\gamma})\mathbf{Q}_n^o(\boldsymbol{\gamma}))^{-1} \\
&\quad\times \mathbf{Q}_n^o(\boldsymbol{\gamma})'\mathbf{P}_n^o(\boldsymbol{\gamma})\mathbf{V}_n^o(\boldsymbol{\gamma})\mathbf{P}_n^o(\boldsymbol{\gamma})\mathbf{Q}_n^o(\boldsymbol{\gamma})\,(\mathbf{Q}_n^o(\boldsymbol{\gamma})'\mathbf{P}_n^o(\boldsymbol{\gamma})\mathbf{Q}_n^o(\boldsymbol{\gamma}))^{-1}
\end{align*}
\emph{and $\mathbf{D}_n^o(\boldsymbol{\gamma})^{-1}$ are $O(1)$.}
\end{theorem}

\begin{proof}
Apply the proof of Exercise~4.26 for each $P^o$ in $\mathcal{P}$ and $\boldsymbol{\gamma}$ in $\Gamma$.
\end{proof}

This result is stated with features that permit fairly general applications,
discussed further next. In particular, we let $\{\mathbf{Z}_{nt}^o(\boldsymbol{\gamma})\}$ be a double array
of matrices. In some cases a single array (sequence) $\{\mathbf{Z}_t^o(\boldsymbol{\gamma})\}$ suffices. In
such cases we can simply put $\mathbf{m}_n^o(\boldsymbol{\gamma}) = n^{-1}\mathbf{Z}^o(\boldsymbol{\gamma})'\bfeps$, so that $\mathbf{V}_n^o(\boldsymbol{\gamma}) \equiv
\mathrm{var}^o(n^{-1/2}\mathbf{Z}^o(\boldsymbol{\gamma})'\bfeps)$, parallel to Exercise~4.26. In this case we also have
$\mathbf{Q}_n^o(\boldsymbol{\gamma}) = E^o(\mathbf{Z}^o(\boldsymbol{\gamma})'\mathbf{X}^o/n)$. These identifications will suffice for now, but
the additional flexibility permitted in Theorem~4.54 will become useful
shortly. Recall that the proof of Exercise~4.26 gives
\begin{align*}
\sqrt{n}\left(\tilde{\bfbeta}_n(\boldsymbol{\gamma}) - \bfbeta^o\right) &= (\mathbf{Q}_n^o(\boldsymbol{\gamma})'\mathbf{V}_n^o(\boldsymbol{\gamma})^{-1}\mathbf{Q}_n^o(\boldsymbol{\gamma}))^{-1} \\
&\quad\times \mathbf{Q}_n^o(\boldsymbol{\gamma})'\mathbf{V}_n^o(\boldsymbol{\gamma})^{-1}n^{-1/2}\mathbf{Z}^o(\boldsymbol{\gamma})'\bfeps + o_{p^o}(1)
\end{align*}
when we choose $\hat{\mathbf{P}}_n(\boldsymbol{\gamma})$ so that $\mathbf{P}_n^o(\boldsymbol{\gamma}) = \mathbf{V}_n^o(\boldsymbol{\gamma})^{-1}$, the asymptotically efficient choice. Comparing this expression to \eqref{eq:4.1} we see that $\tilde{\bfbeta}_n(\boldsymbol{\gamma})$ corresponds to $\hat{\boldsymbol{\theta}}_n(\boldsymbol{\gamma})$, $\mathcal{E}$ is the set of estimators $\mathcal{E} \equiv \{\{\tilde{\bfbeta}_n(\boldsymbol{\gamma})\},\, \boldsymbol{\gamma} \in \Gamma\}$, $\bfbeta^o$
corresponds to $\boldsymbol{\theta}^o$, and we can put
\begin{eqnarray*}
\mathbf{H}_n^o(\boldsymbol{\gamma}) &=& \mathbf{Q}_n^o(\boldsymbol{\gamma})'\mathbf{V}_n^o(\boldsymbol{\gamma})^{-1}\mathbf{Q}_n^o(\boldsymbol{\gamma})\\
\mathbf{s}_n^o(\boldsymbol{\gamma}) &=& \mathbf{Q}_n^o(\boldsymbol{\gamma})'\mathbf{V}_n^o(\boldsymbol{\gamma})^{-1}n^{-1}\mathbf{Z}^o(\boldsymbol{\gamma})'\bfeps\\
\mathbf{I}_n^o(\boldsymbol{\gamma}) &=& \mathbf{H}_n^o(\boldsymbol{\gamma}).
\end{eqnarray*}
Our consideration of the construction of the optimal instrumental variables
estimator makes essential use of these correspondences.

The key to applying the Bates--White constructive method here is the
following result, a version of Theorem~2.6 of Bates and White (1993).

\begin{proposition}\label{prop:4.55}
Let $\hat{\boldsymbol{\theta}}_n(\boldsymbol{\gamma})$ satisfy \eqref{eq:4.1} for all $\boldsymbol{\gamma}$ in a nonempty set $\Gamma$ and all $P^o$ in $\mathcal{P}$, and suppose there exists $\boldsymbol{\gamma}^*$ in $\Gamma$ such that
\begin{equation}\label{eq:4.2}
\mathbf{H}_n^o(\boldsymbol{\gamma}) = \mathrm{cov}^o\left(\sqrt{n}\mathbf{s}_n^o(\boldsymbol{\gamma}), \sqrt{n}\mathbf{s}_n^o(\boldsymbol{\gamma}^*)\right)
\end{equation}
for all $\boldsymbol{\gamma} \in \Gamma$ and all $P^o$ in $\mathcal{P}$. Then for all $\boldsymbol{\gamma} \in \Gamma$ and all $P^o$ in $\mathcal{P}$
\[
\mathrm{avar}^o\left(\hat{\boldsymbol{\theta}}_n(\boldsymbol{\gamma})\right) - \mathrm{avar}^o\left(\hat{\boldsymbol{\theta}}_n(\boldsymbol{\gamma}^*)\right)
\]
is positive semidefinite for all $n$ sufficiently large.
\end{proposition}

\begin{proof}
Given \eqref{eq:4.1}, we have that for all $\boldsymbol{\gamma}$ in $\Gamma$, all $P^o$ in $\mathcal{P}$, and all $n$
sufficiently large
\begin{eqnarray*}
\mathrm{avar}^o\left(\hat{\boldsymbol{\theta}}_n(\boldsymbol{\gamma})\right) &=& \mathbf{H}_n^o(\boldsymbol{\gamma})^{-1}\mathbf{I}_n^o(\boldsymbol{\gamma})\mathbf{H}_n^o(\boldsymbol{\gamma})^{-1}\\
\mathrm{avar}^o\left(\hat{\boldsymbol{\theta}}_n(\boldsymbol{\gamma}^*)\right) &=& \mathbf{H}_n^o(\boldsymbol{\gamma}^*)^{-1}\mathbf{I}_n^o(\boldsymbol{\gamma}^*)\mathbf{H}_n^o(\boldsymbol{\gamma}^*)^{-1} = \mathbf{H}_n^o(\boldsymbol{\gamma}^*)^{-1},
\end{eqnarray*}
so
\[
\mathrm{avar}^o\left(\hat{\boldsymbol{\theta}}_n(\boldsymbol{\gamma})\right) - \mathrm{avar}^o\left(\hat{\boldsymbol{\theta}}_n(\boldsymbol{\gamma}^*)\right) = \mathbf{H}_n^o(\boldsymbol{\gamma})^{-1}\mathbf{I}_n^o(\boldsymbol{\gamma})\mathbf{H}_n^o(\boldsymbol{\gamma})^{-1} - \mathbf{H}_n^o(\boldsymbol{\gamma}^*)^{-1}.
\]
We now show that this quantity equals
\[
\mathrm{var}^o(\mathbf{H}_n^o(\boldsymbol{\gamma})^{-1}\sqrt{n}\mathbf{s}_n^o(\boldsymbol{\gamma}) - \mathbf{H}_n^o(\boldsymbol{\gamma}^*)^{-1}\sqrt{n}\mathbf{s}_n^o(\boldsymbol{\gamma}^*)).
\]
Note that \eqref{eq:4.2} implies
\[
\mathrm{var}^o\left(\begin{array}{c}\sqrt{n}\mathbf{s}_n^o(\boldsymbol{\gamma})\\\sqrt{n}\mathbf{s}_n^o(\boldsymbol{\gamma}^*)\end{array}\right) = \begin{pmatrix} \mathbf{I}_n^o(\boldsymbol{\gamma}) & \mathbf{H}_n^o(\boldsymbol{\gamma}) \\ \mathbf{H}_n^o(\boldsymbol{\gamma})' & \mathbf{H}_n^o(\boldsymbol{\gamma}^*) \end{pmatrix},
\]
So the variance of $\mathbf{H}_n^o(\boldsymbol{\gamma})^{-1}\sqrt{n}\mathbf{s}_n^o(\boldsymbol{\gamma}) - \mathbf{H}_n^o(\boldsymbol{\gamma}^*)^{-1}\sqrt{n}\mathbf{s}_n^o(\boldsymbol{\gamma}^*)$ is given by
\begin{align*}
&\mathbf{H}_n^o(\boldsymbol{\gamma})^{-1}\mathbf{I}_n^o(\boldsymbol{\gamma})\mathbf{H}_n^o(\boldsymbol{\gamma})^{-1} + \mathbf{H}_n^o(\boldsymbol{\gamma}^*)^{-1}\mathbf{H}_n^o(\boldsymbol{\gamma}^*)\mathbf{H}_n^o(\boldsymbol{\gamma}^*)^{-1} \\
&\quad - \mathbf{H}_n^o(\boldsymbol{\gamma})^{-1}\mathbf{H}_n^o(\boldsymbol{\gamma})\mathbf{H}_n^o(\boldsymbol{\gamma}^*)^{-1} - \mathbf{H}_n^o(\boldsymbol{\gamma}^*)^{-1}\mathbf{H}_n^o(\boldsymbol{\gamma})'\mathbf{H}_n^o(\boldsymbol{\gamma})^{-1} \\
&\qquad = \mathbf{H}_n^o(\boldsymbol{\gamma})^{-1}\mathbf{I}_n^o(\boldsymbol{\gamma})\mathbf{H}_n^o(\boldsymbol{\gamma})^{-1} - \mathbf{H}_n^o(\boldsymbol{\gamma}^*)^{-1},
\end{align*}
as claimed. Since this quantity is a variance-covariance matrix it is positive
semidefinite, and the proof is complete.
\end{proof}

The key to finding the asymptotically efficient estimator in a given class
is thus to find $\boldsymbol{\gamma}^*$ such that for all $P^o$ in $\mathcal{P}$ and $\boldsymbol{\gamma}$ in $\Gamma$
\[
\mathbf{H}_n^o(\boldsymbol{\gamma}) = \mathrm{cov}^o\left(\sqrt{n}\mathbf{s}_n^o(\boldsymbol{\gamma}), \sqrt{n}\mathbf{s}_n^o(\boldsymbol{\gamma}^*)\right).
\]

Finding the optimal instrumental variables turns out to be somewhat
involved, but key insight can be gained by initially considering $\mathbf{Z}_t^o(\boldsymbol{\gamma})$ as
previously suggested and by requiring that the errors $\{\eps_t\}$ be sufficiently
well behaved. Thus, consider applying \eqref{eq:4.2} with $\mathbf{Z}_t^o(\boldsymbol{\gamma})$. The correspondences previously identified give
\begin{align*}
\mathbf{Q}_n^o(\boldsymbol{\gamma})'\mathbf{V}_n^o(\boldsymbol{\gamma})^{-1}\mathbf{Q}_n^o(\boldsymbol{\gamma}) &= \\
\mathrm{cov}^o\Big(\mathbf{Q}_n^o(\boldsymbol{\gamma})'\mathbf{V}_n^o(\boldsymbol{\gamma})^{-1}n^{-1/2}\mathbf{Z}^o(\boldsymbol{\gamma})'\bfeps,\, &\mathbf{Q}_n^o(\boldsymbol{\gamma}^*)'\mathbf{V}_n^o(\boldsymbol{\gamma}^*)^{-1}n^{-1/2}\mathbf{Z}^o(\boldsymbol{\gamma}^*)'\bfeps\Big).
\end{align*}
Defining $\mathbf{C}_n^o(\boldsymbol{\gamma}, \boldsymbol{\gamma}^*) \equiv \mathrm{cov}^o(n^{-1/2}\mathbf{Z}^o(\boldsymbol{\gamma})'\bfeps, n^{-1/2}\mathbf{Z}^o(\boldsymbol{\gamma}^*)'\bfeps)$, we have the simpler expression
\[
\mathbf{Q}_n^o(\boldsymbol{\gamma})'\mathbf{V}_n^o(\boldsymbol{\gamma})^{-1}\mathbf{Q}_n^o(\boldsymbol{\gamma}) = \mathbf{Q}_n^o(\boldsymbol{\gamma})'\mathbf{V}_n^o(\boldsymbol{\gamma})^{-1}\mathbf{C}_n^o(\boldsymbol{\gamma}, \boldsymbol{\gamma}^*)\,\mathbf{V}_n^o(\boldsymbol{\gamma}^*)^{-1}\mathbf{Q}_n^o(\boldsymbol{\gamma}^*).
\]
It therefore suffices to find $\boldsymbol{\gamma}^*$ such that for all $\boldsymbol{\gamma}$
\begin{equation}\label{eq:4.3}
\mathbf{Q}_n^o(\boldsymbol{\gamma}) = \mathbf{C}_n^o(\boldsymbol{\gamma}, \boldsymbol{\gamma}^*)\,\mathbf{V}_n^o(\boldsymbol{\gamma}^*)^{-1}\mathbf{Q}_n^o(\boldsymbol{\gamma}^*).
\end{equation}

To make this choice tractable, we impose plausible structure on $\eps_t$.
Specifically, we suppose there is an increasing sequence of $\sigma$-fields $\{\mathcal{F}_t\}$
adapted to $\{\eps_t\}$ such that for each $P^o$ in $\mathcal{P}$ $\{\eps_t, \mathcal{F}_t\}$ is a martingale difference sequence. That is, $E^o(\eps_t|\mathcal{F}_{t-1}) = 0$. Recall that the instrumental
variables $\{\mathbf{Z}_t^o\}$ must be uncorrelated with $\eps_t$. It is therefore natural to restrict attention to instrumental variables $\mathbf{Z}_t^o(\boldsymbol{\gamma})$ that are measurable-$\mathcal{F}_{t-1}$,
as then
\begin{eqnarray*}
E^o(\mathbf{Z}_t^o(\boldsymbol{\gamma})\eps_t) &=& E^o(E^o(\mathbf{Z}_t^o(\boldsymbol{\gamma})\eps_t|\mathcal{F}_{t-1}))\\
&=& E^o(\mathbf{Z}_t^o(\boldsymbol{\gamma})E^o(\eps_t|\mathcal{F}_{t-1})) = \mathbf{0}.
\end{eqnarray*}

Imposing these requirements on $\{\eps_t\}$ and $\{\mathbf{Z}_t^o(\boldsymbol{\gamma})\}$, we have that
\begin{align*}
\mathbf{C}_n^o(\boldsymbol{\gamma}, \boldsymbol{\gamma}^*) &= \mathrm{cov}^o\left(n^{-1/2}\mathbf{Z}^o(\boldsymbol{\gamma})'\bfeps, n^{-1/2}\mathbf{Z}^o(\boldsymbol{\gamma}^*)'\bfeps\right) \\
&= n^{-1}\sum_{t=1}^{n}E^o\left(\mathbf{Z}_t^o(\boldsymbol{\gamma})\eps_t\eps_t'\mathbf{Z}_t^o(\boldsymbol{\gamma}^*)'\right) \\
&= n^{-1}\sum_{t=1}^{n}E^o\left(E^o(\mathbf{Z}_t^o(\boldsymbol{\gamma})\eps_t\eps_t'\mathbf{Z}_t^o(\boldsymbol{\gamma}^*)'|\mathcal{F}_{t-1})\right) \\
&= n^{-1}\sum_{t=1}^{n}E^o\left(\mathbf{Z}_t^o(\boldsymbol{\gamma})E^o(\eps_t\eps_t'|\mathcal{F}_{t-1})\mathbf{Z}_t^o(\boldsymbol{\gamma}^*)'\right) \\
&= n^{-1}\sum_{t=1}^{n}E^o\left(\mathbf{Z}_t^o(\boldsymbol{\gamma})\,\bfOmega_t^o\,\mathbf{Z}_t^o(\boldsymbol{\gamma}^*)'\right),
\end{align*}
where $\bfOmega_t^o \equiv E^o(\eps_t\eps_t'|\mathcal{F}_{t-1})$. The first equality is by definition, the second
follows from the martingale difference property of $\eps_t$, the third uses the law
of iterated expectations, and the fourth applies Proposition~3.63.

Consequently, we seek $\boldsymbol{\gamma}^*$ such that
\[
n^{-1}\sum_{t=1}^{n}E^o(\mathbf{Z}_t^o(\boldsymbol{\gamma})\mathbf{X}_t^{o\prime}) = n^{-1}\sum_{t=1}^{n}E^o(\mathbf{Z}_t^o(\boldsymbol{\gamma})\,\bfOmega_t^o\,\mathbf{Z}_t^o(\boldsymbol{\gamma}^*)')\mathbf{V}_n^o(\boldsymbol{\gamma}^*)^{-1}\mathbf{Q}_n^o(\boldsymbol{\gamma}^*),
\]
or, by the law of iterated expectations,
\[
n^{-1}\sum_{t=1}^{n}E^o\left(E^o\left(\mathbf{Z}_t^o(\boldsymbol{\gamma})\mathbf{X}_t^{o\prime}|\mathcal{F}_{t-1}\right)\right)
= n^{-1}\sum_{t=1}^{n}E^o\left(\mathbf{Z}_t^o(\boldsymbol{\gamma})E^o\left(\mathbf{X}_t^{o\prime}|\mathcal{F}_{t-1}\right)\right)
= n^{-1}\sum_{t=1}^{n}E^o\left(\mathbf{Z}_t^o(\boldsymbol{\gamma})\tilde{\mathbf{X}}_t^{o\prime}\right),
\]
where we define $\tilde{\mathbf{X}}_t^o \equiv E^o(\mathbf{X}_t^o|\mathcal{F}_{t-1})$.

Inspecting the last equation, we see that equality holds provided we
choose $\boldsymbol{\gamma}^*$ such that
\[
\mathbf{Q}_n^o(\boldsymbol{\gamma}^*)'\mathbf{V}_n^o(\boldsymbol{\gamma}^*)^{-1}\mathbf{Z}_t^o(\boldsymbol{\gamma}^*)\bfOmega_t^o = \tilde{\mathbf{X}}_t^o.
\]
If $\bfOmega_t^o \equiv E^o(\eps_t\eps_t'|\mathcal{F}_{t-1})$ is nonsingular, as is plausible, this becomes
\[
\mathbf{Q}_n^o(\boldsymbol{\gamma}^*)'\mathbf{V}_n^o(\boldsymbol{\gamma}^*)^{-1}\mathbf{Z}_t^o(\boldsymbol{\gamma}^*) = \tilde{\mathbf{X}}_t^o\bfOmega_t^{o\,-1}
\]
which looks promising, as both sides are measurable-$\mathcal{F}_{t-1}$. We must therefore choose $\mathbf{Z}_t^o(\boldsymbol{\gamma}^*)$ so that a particular linear combination of $\mathbf{Z}_t^o(\boldsymbol{\gamma}^*)$ equals
$\tilde{\mathbf{X}}_t^o\bfOmega_t^{o\,-1}$.

This still looks challenging, however, because $\boldsymbol{\gamma}^*$ also appears in $\mathbf{Q}_n^o(\boldsymbol{\gamma}^*)$
and $\mathbf{V}_n^o(\boldsymbol{\gamma}^*)$. Nevertheless, consider the simplest choice for $\mathbf{Z}_t^o(\boldsymbol{\gamma}^*)$,
\[
\mathbf{Z}_t^o(\boldsymbol{\gamma}^*) = \tilde{\mathbf{X}}_t^o\bfOmega_t^{o\,-1}.
\]
For this to be valid, it must happen that $\mathbf{Q}_n^o(\boldsymbol{\gamma}^*)' = \mathbf{V}_n^o(\boldsymbol{\gamma}^*)$, so that
$\mathbf{Q}_n^o(\boldsymbol{\gamma}^*)'\mathbf{V}_n^o(\boldsymbol{\gamma}^*)^{-1} = \mathbf{I}$. Now with this choice for $\mathbf{Z}_t^o(\boldsymbol{\gamma}^*)$

\begin{align*}
\mathbf{V}_n^o(\boldsymbol{\gamma}^*) &= \mathrm{var}^o\left(n^{-1/2}\sum_{t=1}^{n}\mathbf{Z}_t^o(\boldsymbol{\gamma}^*)\eps_t\right) \\
&= n^{-1}\sum_{t=1}^{n}E^o\left(\mathbf{Z}_t^o(\boldsymbol{\gamma}^*)\eps_t\eps_t'\mathbf{Z}_t^o(\boldsymbol{\gamma}^*)'\right) \\
&= n^{-1}\sum_{t=1}^{n}E^o\left(\mathbf{Z}_t^o(\boldsymbol{\gamma}^*)\bfOmega_t^o\mathbf{Z}_t^o(\boldsymbol{\gamma}^*)'\right) \\
&= n^{-1}\sum_{t=1}^{n}E^o\left(\tilde{\mathbf{X}}_t^o\bfOmega_t^{o\,-1}\bfOmega_t^o\bfOmega_t^{o\,-1}\tilde{\mathbf{X}}_t^{o\prime}\right) \\
&= n^{-1}\sum_{t=1}^{n}E^o\left(\tilde{\mathbf{X}}_t^o\bfOmega_t^{o\,-1}\tilde{\mathbf{X}}_t^{o\prime}\right) \\
&= n^{-1}\sum_{t=1}^{n}E^o\left(\mathbf{Z}_t^o(\boldsymbol{\gamma}^*)\tilde{\mathbf{X}}_t^{o\prime}\right) \\
&= n^{-1}\sum_{t=1}^{n}E^o\left(\mathbf{Z}_t^o(\boldsymbol{\gamma}^*)E^o(\mathbf{X}_t^{o\prime}|\mathcal{F}_{t-1})\right) \\
&= n^{-1}\sum_{t=1}^{n}E^o\left(\mathbf{Z}_t^o(\boldsymbol{\gamma}^*)\mathbf{X}_t^{o\prime}\right) \\
&= \mathbf{Q}_n^o(\boldsymbol{\gamma}^*).
\end{align*}
The choice $\mathbf{Z}_t^o(\boldsymbol{\gamma}^*) = \tilde{\mathbf{X}}_t^o\bfOmega_t^{o\,-1}$ thus satisfies the sufficient condition \eqref{eq:4.2},
and therefore delivers the optimal instrumental variables.

Nevertheless, using these optimal instrumental variables presents apparent obstacles in that neither $\bfOmega_t^o \equiv E^o(\eps_t\eps_t'|\mathcal{F}_{t-1})$ nor $\tilde{\mathbf{X}}_t^o = E^o(\mathbf{X}_t^o|\mathcal{F}_{t-1})$
are necessarily known or observable. We shall see that these obstacles can
be overcome under reasonable assumptions. Before dealing with the general
case, however, we consider two important special cases in which these obstacles are removed. This provides insight useful in dealing with the general
case, and leads us naturally to procedures having general applicability.

Thus, suppose for the moment
\[
E^o(\eps_t\eps_t'|\mathcal{F}_{t-1}) = \bfOmega_t, \quad t = 1, 2, \ldots\,,
\]
where $\bfOmega_t$ is a known $p \times p$ matrix. (We consider the consequences of not
knowing $\bfOmega_t$ later.) Although this restricts $\mathcal{P}$, it does not restrict it completely --- different choices for $P^o$ can imply different distributions for $\eps_t$,
even though $\bfOmega_t$ is identical.

If we then also require that $\mathbf{X}_t^o$ is measurable-$\mathcal{F}_{t-1}$ it follows that
\[
\tilde{\mathbf{X}}_t^o = E^o(\mathbf{X}_t^o|\mathcal{F}_{t-1}) = \mathbf{X}_t^o,
\]
so that $\tilde{\mathbf{X}}_t^o = \mathbf{X}_t^o$ is observable. Note that this also implies
\[
E^o(\mathbf{Y}_t^o|\mathcal{F}_{t-1}) = E^o(\mathbf{X}_t^{o\prime}\bfbeta^o + \eps_t|\mathcal{F}_{t-1}) = \mathbf{X}_t^{o\prime}\bfbeta^o,
\]
so that the conditional mean of $\mathbf{Y}_t^o$ still depends on $P^o$ through $\bfbeta^o$, but we
can view $(i')$ as defining a relation useful for prediction. This requirement
rules out a system of simultaneous equations, but still permits $\mathbf{X}_t^o$ to contain
lags of elements of $\mathbf{Y}_t^o$.

With these restrictions on $\mathcal{P}$ we have that $\mathbf{Z}_t^o(\boldsymbol{\gamma}^*) = \tilde{\mathbf{X}}_t^o\bfOmega_t^{o\,-1} = \mathbf{X}_t^o\bfOmega_t^{-1}$
delivers a usable set of optimal instrumental variables. Observe that this
clearly illustrates the way in which the existence of an optimal IV estimator
depends on the collection $\mathcal{P}$ of candidate data generating processes.

Let us now examine the estimator resulting from choosing the optimal
instrumental variables in this setting.

\begin{exercises}
\item\label{ex:4.56} \textit{Let $\mathbf{X}_t$ be $\mathcal{F}_{t-1}$-measurable and suppose that $\bfOmega_t = E^o(\eps_t\eps_t'|\mathcal{F}_{t-1})$ for all $P^o$ in $\mathcal{P}$. Show that with the choice $\mathbf{Z}_t^o(\boldsymbol{\gamma}^*) = \mathbf{X}_t^o\bfOmega_t^{-1}$ then $\mathbf{V}_n(\boldsymbol{\gamma}^*) = E(n^{-1}\sum_{t=1}^{n}\mathbf{X}_t^o\bfOmega_t^{-1}\mathbf{X}_t^{o\prime})$ and show that with $\hat{\mathbf{V}}_n(\boldsymbol{\gamma}^*) = n^{-1}\sum_{t=1}^{n}\mathbf{X}_t^o\bfOmega_t^{-1}\mathbf{X}_t^{o\prime}$ this yields the GLS estimator}
\begin{eqnarray*}
\bfbeta_n^* &=& \left(n^{-1}\sum_{t=1}^{n}\mathbf{X}_t^o\bfOmega_t^{-1}\mathbf{X}_t^{o\prime}\right)^{-1}n^{-1}\sum_{t=1}^{n}\mathbf{X}_t^o\bfOmega_t^{-1}\mathbf{Y}_t^o\\
&=& \left(\mathbf{X}'\bfOmega^{-1}\mathbf{X}\right)^{-1}\mathbf{X}'\bfOmega^{-1}\mathbf{Y},
\end{eqnarray*}
\textit{where $\bfOmega$ is the $np \times np$ block diagonal matrix with $p \times p$ diagonal blocks $\bfOmega_t$, $t = 1, \ldots, n$ and zeroes elsewhere, and we drop the ${}^o$ superscript in writing $\mathbf{X}$ and $\mathbf{Y}$.}
\end{exercises}

This is the generalized least squares estimator whose finite sample efficiency properties we discussed in Chapter~1. Our current discussion shows
that these finite sample properties have large sample analogs in a setting
that neither requires normality for the errors $\eps_t$ nor requires that the regressors $\mathbf{X}_t^o$ be nonstochastic.

In fact, we have now proven part $(a)$ of the following asymptotic efficiency
result for the generalized least squares estimator.

\begin{theorem}[Efficiency of GLS]\label{thm:4.57}
Suppose the conditions of Theorem~4.54 hold and that in addition $\mathcal{P}$ is such that for each $P^o$ in $\mathcal{P}$
\begin{enumerate}[label=(\roman*)]
\item $\mathbf{Y}_t^o = \mathbf{X}_t^{o\prime}\bfbeta^o + \eps_t$, \quad $t = 1, 2, \ldots$\,, $\bfbeta^o \in \mathbb{R}^k$,
where $\{(\mathbf{X}_{t+1}^o, \eps_t), \mathcal{F}_t\}$ is an adapted stochastic sequence with $E^o(\eps_t|
\mathcal{F}_{t-1}) = \mathbf{0}$, $t = 1, 2, \ldots$\,; and
\item $E^o(\eps_t\eps_t'|\mathcal{F}_{t-1}) = \bfOmega_t$ is a known finite and nonsingular $p \times p$ matrix, $t = 1, 2, \ldots$\,.
\item For each $\boldsymbol{\gamma} \in \Gamma$, $\mathbf{Z}_t^o(\boldsymbol{\gamma})$ is $\mathcal{F}_{t-1}$-measurable, $t = 1, 2, \ldots$\,, and $\mathbf{Q}_n^o(\boldsymbol{\gamma}) = n^{-1}\sum_{t=1}^{n}E(\mathbf{Z}_t^o(\boldsymbol{\gamma})\mathbf{X}_t^{o\prime})$.
\end{enumerate}

\emph{Then}
\begin{enumerate}[label=(\alph*)]
\item $\mathbf{Z}_t^o(\boldsymbol{\gamma}^*) = \mathbf{X}_t^o\bfOmega_t^{-1}$ delivers
\[
\bfbeta_n^* = \left(n^{-1}\sum \mathbf{X}_t^o\bfOmega_t^{-1}\mathbf{X}_t^{o\prime}\right)^{-1}n^{-1}\sum \mathbf{X}_t^o\bfOmega_t^{-1}\mathbf{Y}_t^o,
\]
which is asymptotically efficient in $\mathcal{E}$ for $\mathcal{P}$.
\item Further, the asymptotic covariance matrix of $\bfbeta_n^*$ is given by
\[
\mathrm{avar}^o(\bfbeta_n^*) = \left[n^{-1}\sum_{t=1}^{n}E^o\left(\mathbf{X}_t^o\bfOmega_t^{-1}\mathbf{X}_t^{o\prime}\right)\right]^{-1},
\]
which is consistently estimated by
\[
\hat{\mathbf{D}}_n = \left[n^{-1}\sum_{t=1}^{n}\mathbf{X}_t^o\bfOmega_t^{-1}\mathbf{X}_t^{o\prime}\right]^{-1},
\]
i.e., $\hat{\mathbf{D}}_n - \mathrm{avar}^o(\bfbeta_n^*) \xrightarrow{p^o} \mathbf{0}$.
\end{enumerate}
\end{theorem}

We leave the proof of $(b)$ as an exercise:

\begin{exercises}
\item\label{ex:4.58} \textit{Prove Theorem~4.57 (b).}
\end{exercises}

As we remarked in Chapter~1, the GLS estimator typically is not feasible
because the conditional covariances $\{\bfOmega_t\}$ are usually unknown. Nevertheless, we might be willing to model $\bfOmega_t$ based on an assumed data generating
process for $\eps_t\eps_t'$, say
\[
E^o(\eps_t\eps_t'|\mathcal{F}_{t-1}) = \mathbf{h}(\mathbf{W}_t^o, \boldsymbol{\alpha}^o),
\]
where $\mathbf{h}$ is a known function mapping $\mathcal{F}_{t-1}$-measurable random variables
$\mathbf{W}_t^o$ and unknown parameters $\boldsymbol{\alpha}^o$ into $p \times p$ matrices. Note that different
$P^o$'s can now be accommodated by different values for $\boldsymbol{\alpha}^o$, and that $\mathbf{W}_t^o$
may depend on $P^o$.

An example is the ARCH$(q)$ data generating process of Engle (1982) in
which $\mathbf{h}(\mathbf{W}_t^o, \boldsymbol{\alpha}^o) = \alpha_0^o + \alpha_1^o\eps_{t-1}^2 + \cdots + \alpha_q^o\eps_{t-q}^2$ for the scalar ($p=1$) case.
Here, $\mathbf{W}_t^o = (\eps_{t-1}^2, \ldots, \eps_{t-q}^2)$. The popular GARCH(1,1) DGP of Bollerslev
(1986) can be similarly treated by viewing it as an ARCH($\infty$) DGP, in
which case $\mathbf{W}_t^o = (\eps_{t-1}^2, \eps_{t-2}^2, \ldots)$.

If we can obtain an estimator, say $\hat{\boldsymbol{\alpha}}_n$, consistent for $\boldsymbol{\alpha}^o$ regardless of
which $P^o$ in $\mathcal{P}$ generated the data, then we can hope that replacing the
unknown $\bfOmega_t = \mathbf{h}(\mathbf{W}_t^o, \boldsymbol{\alpha}^o)$ with the estimator $\hat{\bfOmega}_{nt} = \mathbf{h}(\mathbf{W}_t^o, \hat{\boldsymbol{\alpha}}_n)$ might
still deliver an efficient estimator of the form
\[
\bfbeta_n^* = \left(n^{-1}\sum_{t=1}^{n}\mathbf{X}_t^o\hat{\bfOmega}_{nt}^{-1}\mathbf{X}_t^{o\prime}\right)^{-1}n^{-1}\sum_{t=1}^{n}\mathbf{X}_t^o\hat{\bfOmega}_{nt}^{-1}\mathbf{Y}_t^o.
\]
Such estimators are called ``feasible'' GLS (FGLS) estimators, because they
can feasibly be computed whereas, in the absence of knowledge of $\bfOmega_t$, GLS
cannot.

Inspecting the FGLS estimator, we see that it is also an IV estimator
with
\begin{eqnarray*}
\mathbf{Z}_{nt}^o &=& \mathbf{X}_t^o\hat{\bfOmega}_{nt}^{-1}\\
\hat{\mathbf{P}}_n &=& n^{-1}\sum \mathbf{X}_t^o\hat{\bfOmega}_{nt}^{-1}\mathbf{X}_t^{o\prime}.
\end{eqnarray*}
Observe that $\mathbf{Z}_{nt}^o = \mathbf{X}_t^o\hat{\bfOmega}_{nt}^{-1}$ is now double subscripted to make explicit its
dependence on both $n$ and $t$. This double subscripting is explicitly allowed
for in Theorem~4.54, and we can now see the usefulness of the flexibility
this affords. To accommodate this choice as well as other similar choices,
we now elaborate our specifications of $\mathbf{Z}_{nt}^o$ and $\boldsymbol{\gamma}$. Specifically, we put
\[
\mathbf{Z}_{nt}^o(\boldsymbol{\gamma}) = \mathbf{Z}_{t2}^o(\gamma_2)\boldsymbol{\zeta}(\mathbf{Z}_{t1}^o(\gamma_1), \gamma_{3n}).
\]
Now $\boldsymbol{\gamma}$ has three components: $\boldsymbol{\gamma} = (\gamma_1, \gamma_2, \gamma_3)$, where $\gamma_3 = \{\gamma_{3n}\}$ is a sequence of $\mathcal{F}$-measurable mappings for all $n = 1, 2, \ldots$\,, and $\mathbf{Z}_{t1}^o(\gamma_1)$ and $\mathbf{Z}_{t2}^o(\gamma_2)$ are $\mathcal{F}_{t-1}$-measurable for all $\gamma_1$ and $\gamma_2$, $t = 1, 2, \ldots$\,; and there exists $\gamma_3^o$ such that
$\boldsymbol{\zeta}(\mathbf{Z}_{t1}^o(\gamma_1), \gamma_{3n})$ is a $p \times p$ matrix, and $\mathbf{Z}_{t1}^o(\gamma_1)$ is an $l \times p$ random matrix.

By putting $\mathbf{Z}_{t1}^o(\gamma_1) = \mathbf{X}_t^o$, $\mathbf{Z}_{t2}^o(\gamma_2) = \mathbf{W}_t^o$, $\gamma_{3n} = \hat{\boldsymbol{\alpha}}_n$, and $\boldsymbol{\zeta} = (\mathbf{h})^{-1}$ we
obtain
\begin{eqnarray*}
\mathbf{Z}_{nt}^o(\boldsymbol{\gamma}) &=& \mathbf{X}_t^o(\mathbf{h}(\mathbf{W}_t^o, \hat{\boldsymbol{\alpha}}_n))^{-1}\\
&=& \mathbf{X}_t^o\hat{\bfOmega}_{nt}^{-1}.
\end{eqnarray*}
Many other choices that yield consistent asymptotically normal IV estimators are also permitted.

The requirement that $\boldsymbol{\zeta}$ be known is not restrictive; in particular, we can
choose $\boldsymbol{\zeta}$ such that $\boldsymbol{\zeta}(\mathbf{z}_1, \gamma_{3n}) = \gamma_{3n}(\mathbf{z}_1)$ (the evaluation functional) which
allows consideration of nonparametric estimators of the inverse covariance
matrix, $\bfOmega_t^{o\,-1}$, as well as parametric estimators such as ARCH (set $\gamma_{3n}(\cdot) =
[h(\cdot, \hat{\boldsymbol{\alpha}}_n)]^{-1}$).

To see what choices for $\boldsymbol{\gamma}$ are permitted, it suffices to apply Theorem
4.54 with $\mathbf{Z}^o(\boldsymbol{\gamma})$ interpreted as the $np \times l$ matrix with blocks $\mathbf{Z}_{nt}^o(\boldsymbol{\gamma})'$. We
expand $\mathcal{P}$ to include data generating processes such that $\bfOmega_t^{o\,-1}$ can be
sufficiently well estimated by $\boldsymbol{\zeta}(\mathbf{Z}_{t1}^o(\gamma_1), \gamma_{3n})$ for some choice of $(\gamma_1, \gamma_3)$.
(What is meant by ``sufficiently well'' is that conditions $(ii)$ and $(iii)$
of Theorem~4.54 hold.) It is here that the generality afforded by Theorem~4.54 $(ii)$ is required, as we can no longer necessarily take $\mathbf{m}_n^o(\boldsymbol{\gamma}) =
n^{-1}\sum_{t=1}^{n}\mathbf{Z}_{nt}^o(\boldsymbol{\gamma})\eps_t$. The presence of $\gamma_{3n}$ and the nature of the function
$\boldsymbol{\zeta}$ may cause $n^{-1/2}\sum_{t=1}^{n}\mathbf{Z}_{nt}^o(\boldsymbol{\gamma})\eps_t$ to fail to have a finite second moment,
even though $n^{-1/2}\sum_{t=1}^{n}\mathbf{Z}_{nt}^o(\boldsymbol{\gamma})\eps_t - n^{1/2}\mathbf{m}_n^o(\boldsymbol{\gamma}) = o_{p^o}(1)$ for well behaved
$\mathbf{m}_n^o(\boldsymbol{\gamma})$. Similarly, there is no guarantee that $E^o(\mathbf{Z}_{nt}^o(\boldsymbol{\gamma})\mathbf{X}_t^{o\prime})$ exists, although the sample average $n^{-1}\sum_{t=1}^{n}\mathbf{Z}_{nt}^o(\boldsymbol{\gamma})\mathbf{X}_t^{o\prime}$ may have a well defined
probability limit $\mathbf{Q}_n^o(\boldsymbol{\gamma})$.

Consider first the form of $\mathbf{m}_n^o$. For notational economy, let $\gamma_{t1}^o \equiv \mathbf{Z}_{t1}^o(\gamma_1)$
and $\gamma_{t2}^o \equiv \mathbf{Z}_{t2}^o(\gamma_2)$. Then with
$\mathbf{Z}_{nt}^o(\boldsymbol{\gamma}) = \gamma_{t2}^o\boldsymbol{\zeta}(\gamma_{t1}^o, \gamma_{3n})$ we have
\begin{align*}
n^{-1/2}\sum_{t=1}^{n}\mathbf{Z}_{nt}^o(\boldsymbol{\gamma})\eps_t &= n^{-1/2}\sum_{t=1}^{n}\gamma_{t2}^o\boldsymbol{\zeta}(\gamma_{t1}^o, \gamma_{3n})\eps_t\\
&= n^{-1/2}\sum \gamma_{t2}^o\boldsymbol{\zeta}(\gamma_{t1}^o, \gamma_3^o)\eps_t\\
&\quad + n^{-1/2}\sum \gamma_{t2}^o[\boldsymbol{\zeta}(\gamma_{t1}^o, \gamma_{3n}) - \boldsymbol{\zeta}(\gamma_{t1}^o, \gamma_3^o)]\eps_t,
\end{align*}
where we view $\gamma_3^o$ as the probability limit of $\gamma_{3n}$ under $P^o$. If sufficient
regularity is available to ensure that the second term vanishes in probability,
i.e.,
\[
n^{-1/2}\sum_{t=1}^{n}\gamma_{t2}^o[\boldsymbol{\zeta}(\gamma_{t1}^o, \gamma_{3n}) - \boldsymbol{\zeta}(\gamma_{t1}^o, \gamma_3^o)]\eps_t = o_{p^o}(1),
\]
then we can take
\[
n^{1/2}\mathbf{m}_n^o(\boldsymbol{\gamma}) = n^{-1/2}\sum_{t=1}^{n}\gamma_{t2}^o\boldsymbol{\zeta}(\gamma_{t1}^o, \gamma_3^o)\eps_t.
\]
We ensure ``sufficient regularity'' simply by assuming that $\Gamma$ and $\mathcal{P}$ are
chosen so that the desired convergence holds for all $\boldsymbol{\gamma} \in \Gamma$ and $P^o$ in $\mathcal{P}$.

The form of $\mathbf{Q}_n^o$ follows similarly. With $\mathbf{Z}_{nt}^o(\boldsymbol{\gamma}) = \gamma_{t2}^o\boldsymbol{\zeta}(\gamma_{t1}^o, \gamma_{3n})$ we have
\begin{align*}
n^{-1}\sum_{t=1}^{n}\mathbf{Z}_{nt}^o(\boldsymbol{\gamma})\mathbf{X}_t^{o\prime} &= n^{-1}\sum_{t=1}^{n}\gamma_{t2}^o\boldsymbol{\zeta}(\gamma_{t1}^o, \gamma_{3n})\mathbf{X}_t^{o\prime}\\
&= n^{-1}\sum \gamma_{t2}^o\boldsymbol{\zeta}(\gamma_{t1}^o, \gamma_3^o)\mathbf{X}_t^{o\prime}\\
&\quad + n^{-1}\sum [\gamma_{t2}^o\boldsymbol{\zeta}(\gamma_{t1}^o, \gamma_{3n}) - \gamma_{t2}^o\boldsymbol{\zeta}(\gamma_{t1}^o, \gamma_3^o)]\mathbf{X}_t^{o\prime}.
\end{align*}
Provided the second term vanishes in probability ($\text{-}P^o$) and a law of large
numbers applies to the first term,
\[
\mathbf{Q}_n^o(\boldsymbol{\gamma}) = n^{-1}\sum_{t=1}^{n} E^o[\gamma_{t2}^o\boldsymbol{\zeta}(\gamma_{t1}^o, \gamma_3^o)\mathbf{X}_t^{o\prime}].
\]

With $\mathbf{V}_n^o(\boldsymbol{\gamma}) \equiv \mathrm{var}^o(n^{1/2}\mathbf{m}_n^o(\boldsymbol{\gamma}))$, the Bates--White correspondences are
now
\begin{eqnarray*}
\mathbf{H}_n^o(\boldsymbol{\gamma}) &=& \mathbf{Q}_n^o(\boldsymbol{\gamma})'\mathbf{V}_n^o(\boldsymbol{\gamma})^{-1}\mathbf{Q}_n^o(\boldsymbol{\gamma})\\
\mathbf{s}_n^o(\boldsymbol{\gamma}) &=& \mathbf{Q}_n^o(\boldsymbol{\gamma})'\mathbf{V}_n^o(\boldsymbol{\gamma})^{-1}\mathbf{m}_n^o(\boldsymbol{\gamma})\\
\mathbf{I}_n^o(\boldsymbol{\gamma}) &=& \mathbf{H}_n^o(\boldsymbol{\gamma}).
\end{eqnarray*}
(As before, we set $\mathbf{P}_n^o = \mathbf{V}_n^{o\,-1}$.) By argument parallel to that following
Proposition~4.55, the sufficient condition for efficiency is
\[
\mathbf{Q}_n^o(\boldsymbol{\gamma}) = \mathbf{C}_n^o(\boldsymbol{\gamma}, \boldsymbol{\gamma}^*)\,\mathbf{V}_n^o(\boldsymbol{\gamma}^*)^{-1}\mathbf{Q}_n^o(\boldsymbol{\gamma}^*),
\]
where now $\mathbf{C}_n^o(\boldsymbol{\gamma}, \boldsymbol{\gamma}^*) = \mathrm{cov}^o(n^{1/2}\mathbf{m}_n^o(\boldsymbol{\gamma}), n^{1/2}\mathbf{m}_n^o(\boldsymbol{\gamma}^*))$.

Following arguments exactly parallel to those leading to Equation~\eqref{eq:4.3},
we find that the sufficient condition holds if
\[
\mathbf{Q}_n^o(\boldsymbol{\gamma}^*)'\mathbf{V}_n^o(\boldsymbol{\gamma}^*)^{-1}\gamma_3^{*o}(\mathbf{Z}_{t1}^o(\gamma_1^*))\bfOmega_t^o = \mathbf{X}_t^o.
\]
Choosing $\gamma_3^{*o}(\mathbf{Z}_{t1}^o(\gamma_1^*)) = \bfOmega_t^{o\,-1}$ can be verified to ensure
$\mathbf{Q}_n^o(\boldsymbol{\gamma}^*)' = \mathbf{V}_n^o(\boldsymbol{\gamma}^*)$, so optimality follows.

We formalize this result as follows.

\begin{theorem}[Efficiency of FGLS]\label{thm:4.59}
Suppose the conditions of Theorem~4.54 hold and that in addition $\mathcal{P}$ is such that for each $P^o$ in $\mathcal{P}$
\begin{enumerate}[label=(\roman*)]
\item $\mathbf{Y}_t^o = \mathbf{X}_t^{o\prime}\bfbeta^o + \eps_t$, \quad $t = 1, 2, \ldots$\,, $\bfbeta^o \in \mathbb{R}^k$,
where $\{(\mathbf{X}_{t+1}^o, \eps_t), \mathcal{F}_t\}$ is an adapted stochastic sequence with $E^o(\eps_t|
\mathcal{F}_{t-1}) = \mathbf{0}$, $t = 1, 2, \ldots$\,; and
\item $E^o(\eps_t\eps_t'|\mathcal{F}_{t-1}) = \bfOmega_t^o$ is an unknown finite and nonsingular $p \times p$ matrix, $t = 1, 2, \ldots$\,.
\item For each $\boldsymbol{\gamma} \in \Gamma$,
\[
\mathbf{Z}_{nt}^o(\boldsymbol{\gamma}) = \mathbf{Z}_{t2}^o(\gamma_2)\gamma_{3n}(\mathbf{Z}_{t1}^o(\gamma_1)),
\]
where $\boldsymbol{\gamma} \equiv (\gamma_1, \gamma_2, \gamma_3)$, and $\gamma_3 \equiv \{\gamma_{3n}\}$ is such that $\gamma_{3n}$ is $\mathcal{F}$-measurable for all $n = 1, 2, \ldots$\,, and $\mathbf{Z}_{t1}^o(\gamma_1)$ and $\mathbf{Z}_{t2}^o(\gamma_2)$ are $\mathcal{F}_{t-1}$-measurable for all $\gamma_1$ and $\gamma_2$, $t = 1, 2, \ldots$\,; and there exists $\gamma_3^o$ such
that
\[
\mathbf{m}_n^o(\boldsymbol{\gamma}) = n^{-1}\sum \mathbf{Z}_{t2}^o(\gamma_2)\gamma_3^o(\mathbf{Z}_{t1}^o(\gamma_1))\eps_t
\]
and
\[
\mathbf{Q}_n^o(\boldsymbol{\gamma}) = n^{-1}\sum E(\mathbf{Z}_{t2}^o(\gamma_2)\gamma_3^o(\mathbf{Z}_{t1}^o(\gamma_1))\mathbf{X}_t^{o\prime});
\]
\item for some $\gamma_1^*$, $\gamma_3^*$, $\hat{\bfOmega}_{nt}^{-1} = \gamma_{3n}^*(\mathbf{Z}_{t1}^o(\gamma_1^*))$, where $\gamma_3^{*o}(\mathbf{Z}_{t1}^o(\gamma_1^*)) = \bfOmega_t^{o\,-1}$.
\end{enumerate}

\emph{Then}
\begin{enumerate}[label=(\alph*)]
\item $\mathbf{Z}_{nt}(\boldsymbol{\gamma}^*) = \mathbf{X}_t^o\hat{\bfOmega}_{nt}^{-1}$ delivers
\[
\bfbeta_n^* = \left(n^{-1}\sum_{t=1}^{n}\mathbf{X}_t^o\hat{\bfOmega}_{nt}^{-1}\mathbf{X}_t^{o\prime}\right)^{-1}n^{-1}\sum_{t=1}^{n}\mathbf{X}_t^o\hat{\bfOmega}_{nt}^{-1}\mathbf{Y}_t^o,
\]
which is asymptotically efficient in $\mathcal{E}$ for $\mathcal{P}$;
\item the asymptotic covariance matrix of $\bfbeta_n^*$ is given by
\[
\mathrm{avar}^o(\bfbeta_n^*) = \left[n^{-1}\sum_{t=1}^{n}E^o\left(\mathbf{X}_t^o\bfOmega_t^{o\,-1}\mathbf{X}_t^{o\prime}\right)\right]^{-1},
\]
which is consistently estimated by
\[
\hat{\mathbf{D}}_n = \left[n^{-1}\sum_{t=1}^{n}\mathbf{X}_t^o\hat{\bfOmega}_{nt}^{-1}\mathbf{X}_t^{o\prime}\right]^{-1},
\]
i.e., $\hat{\mathbf{D}}_n - \mathrm{avar}^o(\bfbeta_n^*) \xrightarrow{p^o} \mathbf{0}$.
\end{enumerate}
\end{theorem}

\begin{proof}
$(a)$ Apply Proposition~4.55 and the argument preceding the statement of Theorem~4.59. We leave the proof of part $(b)$ as an exercise.
\end{proof}

\begin{exercises}
\item\label{ex:4.60} \textit{Prove part (b) of Theorem~4.59.}

\item\label{ex:4.61} \textit{Suppose conditions 4.59 (i) and 4.59 (ii) hold with $p = 1$ and that}
\[
\bfOmega_t^o = \alpha_1^o + \alpha_2^oW_t^o, \quad t = 1, 2, \ldots\,,
\]
\textit{where $W_t^o = 1$ during recessions and $W_t^o = 0$ otherwise.}
\begin{enumerate}[label=(\roman*)]
\item \textit{Let $\hat{\eps}_t = Y_t^o - \mathbf{X}_t^{o\prime}\hat{\bfbeta}_n$, where $\hat{\bfbeta}_n$ is the OLS estimator, and let $\hat{\boldsymbol{\alpha}}_n = (\hat{\alpha}_{1n}, \hat{\alpha}_{2n})'$ be the OLS estimator from the regression of $\hat{\eps}_t^2$ on a constant and $W_t^o$. Prove that $\hat{\boldsymbol{\alpha}}_n = \boldsymbol{\alpha}^o + o_{p^o}(1)$, where $\boldsymbol{\alpha}^o = (\alpha_1^o, \alpha_2^o)'$, providing any needed additional conditions.}
\item \textit{Next, verify that conditions 4.59 (iii) and (iv) hold for}
\[
\mathbf{Z}_{nt}^o(\boldsymbol{\gamma}^*) = \mathbf{X}_t^o(\hat{\alpha}_{1n} + \hat{\alpha}_{2n}W_t^o)^{-1}.
\]
\end{enumerate}
\end{exercises}

When less is known about the structure of $\bfOmega_t^o$, nonparametric estimators
can be used. See, for example, White and Stinchcombe (1991). The notion
of stochastic equicontinuity (see Andrews, 1994a,b) plays a key role in such
situations, but will not be explored here.

We are now in a position to see how to treat the general situation in
which $\tilde{\mathbf{X}}_t^o \equiv E(\mathbf{X}_t^o|\mathcal{F}_{t-1})$ appears instead of $\mathbf{X}_t^o$. By taking
\[
\mathbf{Z}_{nt}^o(\boldsymbol{\gamma}) = \underset{l \times p}{\gamma_{4n}(\mathbf{Z}_{t2}^o(\gamma_2))}\underset{p \times p}{\gamma_{3n}(\mathbf{Z}_{t1}^o(\gamma_1))}
\]
we can accommodate the choice
\[
\mathbf{Z}_{nt}(\boldsymbol{\gamma}) = \hat{\mathbf{X}}_t\hat{\bfOmega}_{nt}^{-1},
\]
where $\hat{\mathbf{X}}_{nt} = \gamma_{4n}(\mathbf{Z}_{t2}^o(\gamma_2))$ is an estimator of $\tilde{\mathbf{X}}_t^o$. As we might
now expect, this choice delivers the optimal IV estimator.

We can now state our asymptotic efficiency result for the IV estimator.

\begin{theorem}[Efficient IV]\label{thm:4.62}
Suppose the conditions of Theorem~4.54 hold and that in addition $\mathcal{P}$ is such that for each $P^o$ in $\mathcal{P}$
\begin{enumerate}[label=(\roman*)]
\item $\mathbf{Y}_t^o = \mathbf{X}_t^{o\prime}\bfbeta^o + \eps_t$, \quad $t = 1, 2, \ldots$\,, $\bfbeta^o \in \mathbb{R}^k$,
where $\{\eps_t, \mathcal{F}_t\}$ is an adapted stochastic sequence with $E^o(\eps_t|\mathcal{F}_{t-1}) = \mathbf{0}$, $t = 1, 2, \ldots$\,; and
\item $E^o(\eps_t\eps_t'|\mathcal{F}_{t-1}) = \bfOmega_t^o$ is an unknown finite and nonsingular $p \times p$ matrix, $t = 1, 2, \ldots$\,.
\item For each $\boldsymbol{\gamma} \in \Gamma$,
\[
\mathbf{Z}_{nt}^o(\boldsymbol{\gamma}) = \gamma_{4n}(\mathbf{Z}_{t2}^o(\gamma_2))\gamma_{3n}(\mathbf{Z}_{t1}^o(\gamma_1)),
\]
where $\boldsymbol{\gamma} \equiv (\gamma_1, \gamma_2, \gamma_3, \gamma_4)$, and $\gamma_3 \equiv \{\gamma_{3n}\}$ and $\gamma_4 \equiv \{\gamma_{4n}\}$, are such
that $\gamma_{3n}$ and $\gamma_{4n}$ are $\mathcal{F}$-measurable for all $n = 1, 2, \ldots$ and $\mathbf{Z}_{t1}^o(\gamma_1)$
and $\mathbf{Z}_{t2}^o(\gamma_2)$ are $\mathcal{F}_{t-1}$-measurable for all $\gamma_1$ and $\gamma_2$, $t = 1, 2, \ldots$\,; and
there exist $\gamma_3^o$ and $\gamma_4^o$ such that
\[
\mathbf{m}_n^o(\boldsymbol{\gamma}) = n^{-1}\sum \gamma_4^o(\gamma_{t2}^o)\gamma_3^o(\gamma_{t1}^o)\eps_t
\]
and
\[
\mathbf{Q}_n^o(\boldsymbol{\gamma}) = n^{-1}\sum E^o(\gamma_4^o(\gamma_{t2}^o)\gamma_3^o(\gamma_{t1}^o)\mathbf{X}_t^{o\prime});
\]
\item for some $\gamma_1^*$ and $\gamma_3^*$, $\hat{\bfOmega}_{nt}^{-1} = \gamma_{3n}^*(\mathbf{Z}_{t1}^o(\gamma_1^*))$, where $\gamma_3^{*o}(\mathbf{Z}_{t1}^o(\gamma_1^*)) = \bfOmega_t^{o\,-1}$ and for some $\gamma_2^*$ and $\gamma_4^*$, $\hat{\mathbf{X}}_t = \gamma_{4n}^*(\mathbf{Z}_{t2}^o(\gamma_2^*))$, where $\gamma_4^{*o}(\mathbf{Z}_{t2}^o(\gamma_2^*)) = E^o(\mathbf{X}_t^o|\mathcal{F}_{t-1})$.
\end{enumerate}

\emph{Then}
\begin{enumerate}[label=(\alph*)]
\item $\mathbf{Z}_{nt}^o(\boldsymbol{\gamma}^*) = \hat{\mathbf{X}}_{nt}\hat{\bfOmega}_{nt}^{-1}$ delivers
\[
\bfbeta_n^* = \left(n^{-1}\sum_{t=1}^{n}\hat{\mathbf{X}}_{nt}\hat{\bfOmega}_{nt}^{-1}\mathbf{X}_t^{o\prime}\right)^{-1}n^{-1}\sum_{t=1}^{n}\hat{\mathbf{X}}_{nt}\hat{\bfOmega}_{nt}^{-1}\mathbf{Y}_t^o,
\]
which is asymptotically efficient in $\mathcal{E}$ for $\mathcal{P}$;
\item the asymptotic covariance matrix of $\bfbeta_n^*$ is given by
\[
\mathrm{avar}^o(\bfbeta_n^*) = \left[n^{-1}\sum_{t=1}^{n}E^o\left(\tilde{\mathbf{X}}_t^o\bfOmega_t^{o\,-1}\tilde{\mathbf{X}}_t^{o\prime}\right)\right]^{-1},
\]
which is consistently estimated by
\[
\hat{\mathbf{D}}_n = 2\left[n^{-1}\sum \hat{\mathbf{X}}_{nt}\hat{\bfOmega}_{nt}^{-1}\mathbf{X}_t^{o\prime} + n^{-1}\sum \mathbf{X}_t^o\hat{\bfOmega}_{nt}^{-1}\hat{\mathbf{X}}_{nt}'\right]^{-1},
\]
i.e., $\hat{\mathbf{D}}_n - \mathrm{avar}^o(\bfbeta_n^*) \xrightarrow{p^o} \mathbf{0}$.
\end{enumerate}
\end{theorem}

\begin{proof}
$(a)$ Apply Proposition~4.55 and the argument preceding the statement of Theorem~4.62. We leave the proof of part $(b)$ as an exercise.
\end{proof}

\begin{exercises}
\item\label{ex:4.63} \textit{Prove Theorem~4.62 (b).}
\end{exercises}

Other consistent estimators for $\mathrm{avar}^o(\bfbeta_n^*)$ are available, for example
\[
\tilde{\mathbf{D}}_n = \left[n^{-1}\sum_{t=1}^{n}\hat{\mathbf{X}}_{nt}\hat{\bfOmega}_{nt}^{-1}\hat{\mathbf{X}}_{nt}'\right]^{-1}.
\]
This requires additional conditions, however. The estimator $\hat{\mathbf{D}}_n$ above is
consistent without further assumptions.

\begin{exercises}
\item\label{ex:4.64} \textit{Suppose condition 4.62 (i) and (ii) hold and that}
\begin{eqnarray*}
\bfOmega_t^o &=& \boldsymbol{\Sigma}^o, \\
E^o(\mathbf{X}_t^{o\prime}|\mathcal{F}_{t-1}) &=& \mathbf{Z}_t^{o\prime}\boldsymbol{\pi}^o, \quad t = 1, 2, \ldots\,,
\end{eqnarray*}
\textit{where $\boldsymbol{\pi}^o$ is an $l \times k$ matrix.}
\begin{enumerate}[label=(\roman*)]
\item \begin{enumerate}[label=(\alph*)]
\item \textit{Let $\hat{\boldsymbol{\pi}}_n = (n^{-1}\sum_{t=1}^{n}\mathbf{Z}_t^o\mathbf{Z}_t^{o\prime})^{-1}n^{-1}\sum_{t=1}^{n}\mathbf{Z}_t^o\mathbf{X}_t^{o\prime}$.
Give simple conditions under which $\hat{\boldsymbol{\pi}}_n = \boldsymbol{\pi}^o + o_{p^o}(1)$.}
\item \textit{Let $\hat{\boldsymbol{\Sigma}}_n = n^{-1}\sum_{t=1}^{n}\hat{\eps}_t\hat{\eps}_t'$, where $\hat{\eps}_t = \mathbf{Y}_t^o - \mathbf{X}_t^{o\prime}\tilde{\bfbeta}_n$,}
\[
\tilde{\bfbeta}_n = (\mathbf{X}'\mathbf{Z}(\mathbf{Z}'\mathbf{Z})^{-1}\mathbf{Z}'\mathbf{X})^{-1}\mathbf{X}'\mathbf{Z}(\mathbf{Z}'\mathbf{Z})^{-1}\mathbf{Z}'\mathbf{Y}
\]
\textit{(the 2SLS estimator). Give simple conditions under which $\hat{\boldsymbol{\Sigma}}_n = \boldsymbol{\Sigma}^o + o_{p^o}(1)$.}
\end{enumerate}
\item \textit{Next, verify that conditions 4.62 (iii) and 4.62 (iv) hold for $\mathbf{Z}_{nt}^o(\boldsymbol{\gamma}^*) = \hat{\boldsymbol{\pi}}_n'\mathbf{Z}_t^o\hat{\boldsymbol{\Sigma}}_n^{-1}$. The resulting estimator is the three stage least squares (3SLS) estimator.}
\end{enumerate}
\end{exercises}

\section*{References}

\begin{description}
\item Andrews, A. (1994a). ``Asymptotics for semiparametric econometric-models via stochastic equicontinuity.'' \textit{Econometrica}, 62, 43--72.

\item[\textemdash\textemdash\textemdash] (1994b). ``Empirical process methods in econometrics.'' In \textit{Handbook of Econometrics}, 4. Engle, R.~F. and D.~L. McFadden, eds., pp.~2247--2294. North Holland, Amsterdam.

\item Bartle, R.~G. (1976). \textit{The Elements of Real Analysis}. Wiley, New York.

\item Bates C.~E. and H.~White (1993). ``Determination of estimtors with minimum asymptotic covariance matrices.'' \textit{Econometric Theory}, 9, 633--648.

\item Bollerslev, T. (1986). ``Generalized autoregressive conditional heteroskedasticity.'' \textit{Journal of Econometrics}, 31, 307--327.

\item Chamberlain, G. (1982). ``Multivariate regression models for panel data.'' \textit{Journal of Econometrics}, 18, 5--46.

\item Cragg, J. (1983). ``More efficient estimation in the presence of heteroskedasticity of unknown form.'' \textit{Econometrica}, 51, 751--764.

\item Engle, R.~F. (1981). ``Wald, likelihood ratio, and Lagrange multiplier tests in econometrics.'' In \textit{Handbook of Econometrics}, 2, Z.~Griliches and M.~Intrilligator, eds., North Holland, Amsterdam.

\item[\textemdash\textemdash\textemdash] (1982). ``Autoregressive conditional heteroskedasticity with estimates of the variance of United Kingdom inflation.'' \textit{Econometrica}, 50, 987--1008.

\item Goldberger, A.~S. (1964). \textit{Econometric Theory}. Wiley, New York.

\item Hansen, L.~P. (1982). ``Large Sample Properties of Generalized Method of Moments Estimators.'' \textit{Econometrica}, 50, 1029--1054.

\item Lukacs, E. (1970). \textit{Characteristic Functions}. Griffin, London.

\item[\textemdash\textemdash\textemdash] (1975). \textit{Stochastic Convergence}. Academic Press, New York.

\item Rao, C.~R. (1973). \textit{Linear Statistical Inference and Its Applications}. Wiley, New York.

\item White, H. (1982). ``Instrumental variables regression with independent observations.'' \textit{Econometrica}, 50, 483--500.

\item[\textemdash\textemdash\textemdash] (1994). \textit{Estimation, Inference and Specification Analysis}. Cambridge University Press, New York.

\item[\textemdash\textemdash\textemdash] and M.~Stinchcombe (1991). ``Adaptive efficient weighted least squares with dependent observations.'' In \textit{Directions in Robust Statistics and Diagnostics}, W.~Stahel and S.~Weisberg, eds., pp.~337--364. IMA Volumes in Mathematics and Its Applications. Springer-Verlag, New York.
\end{description}
